\chapter{Schemes over valuation rings and models}
\index{model!scheme over a valuation ring}\index{valuation ring}
\label{chap:models}

The global non-Archimedean theory turns analytic questions on $X^{\mathrm{an}}$
into algebraic questions on models of $X$. This chapter establishes the model
language used throughout the book. Two regimes must be distinguished.
Admissible formal geometry works over rather general complete rank-one valued
fields with nontrivial absolute value. Classical intersection theory and the
original BFJ construction are most transparent over a discrete valuation ring. We state explicitly when
discreteness or noetherianity is used.

Starting from the valuation rings of \cref{chap:valuations}, we introduce
models, vertical divisors, and common dominations, then explain flattening
and the availability of SNC or altered models. Intersection numbers and
relative nefness prepare the model metrics of \cref{chap:metrics}; the
closing dictionary keeps their algebraic and analytic roles distinct.

Throughout, $K$ is a complete non-Archimedean field,
\[
 K^\circ=\{a\in K\mid |a|\leq1\},\qquad
 K^{\circ\circ}=\{a\in K\mid |a|<1\},\qquad
 \widetilde K=K^\circ/K^{\circ\circ}.
\]
The valuation is $v_K=-\log|\cdot|$. In sections involving a uniformizer we
assume that $K$ is discretely valued and use the rescaled absolute value
with $|\pi|=\rme^{-1}$, so $v_K(\pi)=1$.

\section{Algebra over a valuation ring}

\begin{proposition}
\label{prop:valuation-ring-basic}
$K^\circ$ is a valuation ring of $K$, its maximal ideal is
$K^{\circ\circ}$, and its residue field is $\widetilde K$.
\end{proposition}
\begin{proof}
For $a\in K^\times$, either $|a|\leq1$ or $|a^{-1}|<1$, so the valuation-ring
criterion of \Cref{prop:valuation-ring-characterization} applies.
An element of $K^\circ$ is a unit exactly when its absolute value is $1$.
Hence the nonunits form $K^{\circ\circ}$, which is the unique maximal ideal.
The last assertion is the definition of $\widetilde K$.
\end{proof}

The ideal theory of a valuation ring is extremely rigid.  The next two
statements isolate the exact amount of rigidity used in this book: finitely
generated ideals are principal, and noetherianity forces discreteness.  Both
proofs are elementary, and writing them out fixes the arguments that later
chapters use silently.

\begin{lemma}
\label{lem:valuation-ideals-principal}
\index{valuation ring!ideals}
Let $R$ be a valuation ring with fraction field $F$ and valuation $v$.
\begin{enumerate}
\item For nonzero $a,b\in R$, one has $a\mid b$ in $R$ if and only if
  $v(a)\leq v(b)$.  In particular any two elements are comparable under
  divisibility.
\item Every finitely generated ideal of $R$ is principal, generated by any
  element of least value in a finite generating set.
\item The ideals of $R$ are totally ordered by inclusion.
\end{enumerate}
\end{lemma}
\begin{proof}
(1) The quotient $b/a\in F$ satisfies $v(b/a)=v(b)-v(a)$, and it lies in $R$
exactly when this difference is nonnegative.

(2) Let $I=(a_1,\ldots,a_m)$ with all $a_i\ne0$ (discard zero generators;
if all vanish, $I=0$ is principal).  Choose $i_0$ with
$v(a_{i_0})=\min_iv(a_i)$.  By (1), $a_{i_0}$ divides every $a_i$, so
$I=(a_{i_0})$.

(3) Let $I,J$ be ideals with $I\not\subset J$, and choose
$a\in I\setminus J$.  For any nonzero $b\in J$ we cannot have
$b\mid a$, since that would put $a$ in $J$.  Hence $a\mid b$ by (1), so
$b\in(a)\subset I$.  Thus $J\subset I$.
\end{proof}

\begin{proposition}
\label{prop:noetherian-valuation-dvr}
\index{valuation ring!noetherian}
A valuation ring is noetherian if and only if it is a field or a discrete
valuation ring.  Consequently $K^\circ$ is noetherian if and only if the
absolute value of $K$ is trivial or discrete.
\end{proposition}
\begin{proof}
A field is noetherian, and a discrete valuation ring is a principal ideal
domain, hence noetherian: its ideals are $\{0\}$ and the powers of the
maximal ideal.

Conversely, let $R$ be a noetherian valuation ring that is not a field, with
valuation $v$ and maximal ideal $\mathfrak m\ne0$.  By
\Cref{lem:valuation-ideals-principal}, every ideal is finitely generated and
hence principal; write $\mathfrak m=(t)$.

\emph{Step 1: $\bigcap_{n\geq1}(t^n)=0$.}  The intersection
$I=\bigcap_n(t^n)$ is an ideal, hence principal, say $I=(c)$.  If $c\ne0$,
then $c\in(t^{n+1})$ for all $n$, so $c/t\in I=(c)$; write $c/t=rc$ with
$r\in R$.  Then $c(1-rt)=0$.  Since $t\in\mathfrak m$, the element $1-rt$ is
a unit, forcing $c=0$, a contradiction.  Hence $I=0$.

\emph{Step 2: $R$ has rank one.}  Suppose $\mathfrak p$ is a nonzero prime
ideal with $\mathfrak p\subsetneq\mathfrak m$ and pick a nonzero
$x\in\mathfrak p$.  For each $n$, either $x\mid t^n$ or $t^n\mid x$.  The
first case would give $t^n\in(x)\subset\mathfrak p$ and hence
$t\in\mathfrak p$, so $\mathfrak m=(t)\subset\mathfrak p$, contradicting
properness.  Therefore $t^n\mid x$ for every $n$, so
$x\in\bigcap_n(t^n)=0$ by Step 1, a contradiction.  Thus $\mathfrak m$ is
the only nonzero prime.

\emph{Step 3: the value group is $\mathbf Z\,v(t)$.}  First, $v(t)$ is the
least positive value: if $0<v(a)<v(t)$ for some $a$, then
$a\in\mathfrak m=(t)$ would give $v(a)\geq v(t)$.  Next, let $a\in R$ with
$v(a)>0$.  If $v(a)>nv(t)$ for every $n$, then $t^n\mid a$ for every $n$
and Step 1 gives $a=0$.  So there is a largest $n\geq1$ with
$v(a)\geq nv(t)$; then $a/t^n\in R$ and
$0\leq v(a)-nv(t)<v(t)$, which forces $v(a)=nv(t)$ by minimality.  Hence
the value group is infinite cyclic and $R$ is a discrete valuation ring.

For the final assertion, recall from \cref{chap:valuations} that the value
group of $K$ is a subgroup of $\mathbf R$.  It is trivial, discrete, or
dense; the noetherian cases are exactly the first two.
\end{proof}

\begin{lemma}
\label{lem:torsion-free-flat-valuation}
A module over a valuation ring is flat if and only if it is torsion-free.
Consequently, the schematic closure of a subscheme of a generic fibre is flat
over the valuation ring.
\end{lemma}
\begin{proof}
Let $R$ be a valuation ring.  A flat $R$-module is torsion-free because
tensoring the injection $R\xrightarrow{\,a\,}R$ by a flat module preserves
injectivity for every $a\ne0$.

Conversely, let $M$ be torsion-free.  We use the criterion of flatness by
finitely generated ideals: an $R$-module $M$ is flat as soon as the natural
map $I\otimes_RM\to M$ is injective for every finitely generated ideal
$I\subset R$ \cite[Tag 00M5]{stacks-project}.  By
\Cref{lem:valuation-ideals-principal}, every such ideal is principal, say
$I=aR$.  If $a=0$, then $I\otimes_RM=0$ and the criterion is vacuous.  If
$a\ne0$, then $a$ is a nonzerodivisor of the domain $R$, so multiplication
by $a$ is an isomorphism $R\xrightarrow{\ \sim\ }aR$ of $R$-modules.
Tensoring with $M$ identifies $aR\otimes_RM$ with $M$, and under this
identification the map $aR\otimes_RM\to M$ becomes multiplication by $a$ on
$M$.  Torsion-freeness says precisely that this multiplication is
injective.  Thus $M$ is flat.  This is the argument recorded in
\cite[Tag 0539]{stacks-project}.

For the geometric consequence, work on an affine model
$\operatorname{Spec}A$ and let
$I_K\subset A\otimes_RK$ define the generic-fibre subscheme.  Its schematic
closure is defined by the contraction $I=A\cap I_K$, that is, by the kernel
of the composite $A\to A\otimes_RK\to(A\otimes_RK)/I_K$.  The natural map
\[
            A/I\longrightarrow (A\otimes_RK)/I_K
\]
is injective.  The target is a vector space over $K$ and hence is
torsion-free over $R$; so is its submodule $A/I$.  The first part makes
$A/I$ flat over $R$.  Flatness is affine-local, which proves the assertion.
The gluing of the affine constructions into a closed subscheme is carried
out in \Cref{lem:flat-closure-generic-fibre} below.
\end{proof}

\begin{remark}\label{rem:dense-value-group-model-limitations}
If the value group is dense, $K^\circ$ is nonnoetherian. Assertions involving
a uniformizer, multiplicities of the special fibre, or ordinary noetherian intersection
theory are not automatically valid in that regime.
By \Cref{prop:noetherian-valuation-dvr} this is not an occasional pathology
but the general case: for a dense value group $K^\circ$ is \emph{never}
noetherian.  For instance, if $K$ is algebraically closed and nontrivially
valued, its value group is divisible, hence dense.
\Cref{ex:model-dense-value-group} works through a projective model over such
a ring and records exactly which finiteness statements survive.
\end{remark}

\section{Algebraic and formal models}

\begin{definition}\label{def:algebraic-projective-model}
Let $X$ be a proper $K$-scheme. An \emph{algebraic model} of $X$ is a flat,
proper, finitely presented $K^\circ$-scheme $\mathcal X$ together with an
identification $\mathcal X_K\simeq X$. It is a \emph{projective model} if
$\mathcal X\to\operatorname{Spec}K^\circ$ is projective.

If $L$ is a line bundle on $X$, a \emph{model of $(X,L)$} is a pair
$(\mathcal X,\mathcal L)$ in which $\mathcal X$ is a model of $X$ and
$\mathcal L$ is a line bundle on $\mathcal X$ identified with $L$ on the
generic fibre. A model of $mL$ for some $m\geq1$ is called a
\emph{$\mathbf Q$-model} of $L$.
\end{definition}

Affine equations below describe local model charts; properness is imposed
on the global models, including their displayed projective compactifications.

Three comments on this definition are worth fixing once and for all.  First,
the identification $\mathcal X_K\simeq X$ is part of the datum: the same
$K^\circ$-scheme with two different identifications counts as two models,
and morphisms of models are required to respect the identifications.
Second, flatness is not decoration.  Over a valuation ring,
\Cref{lem:torsion-free-flat-valuation} identifies flatness with
torsion-freeness, so requiring it amounts to forbidding components and
torsion sections concentrated in the special fibre; without it the special
fibre could contain arbitrary junk unrelated to $X$, and the special fibre
would not deserve to be called a degeneration of $X$.  Third, over a
discrete valuation ring, finite presentation is the same as finite type,
while over a general $K^\circ$ it is a genuine restriction; that it can
always be achieved is \Cref{prop:flat-closure-finite-presentation} below.

\begin{example}[Flatness deletes vertical junk]
\label{ex:flat-closure-toy}
Let $R=K^\circ$ be a discrete valuation ring and
\[
                 \mathcal V=\operatorname{Spec}R[x]/(\pi x)
                 \subset\mathbf A^1_R.
\]
The element $x$ satisfies $\pi x=0$, so the coordinate ring has
$\pi$-torsion and $\mathcal V$ is not flat.  Geometrically, $\mathcal V$ is
the union of the horizontal section $\{x=0\}\simeq\operatorname{Spec}R$ and
the entire vertical line $\mathbf A^1_{\widetilde K}=\{\pi=0\}$.  Its
generic fibre is the single reduced point $x=0$ of $\mathbf A^1_K$.  The
schematic closure of that generic fibre is cut out by the contraction of
$(x)$, namely $\{x=0\}\simeq\operatorname{Spec}R$: the flat closure has
deleted exactly the vertical component, which is invisible from the generic
fibre.  This tiny example is the mechanism behind
\Cref{lem:torsion-free-flat-valuation} and behind the graph-closure
construction of \Cref{prop:common-domination}.
\end{example}

For the formal generic-fibre constructions in this chapter, assume that
$K$ is nontrivially valued.

\begin{definition}\label{def:admissible-formal-model}
An \emph{admissible formal $K^\circ$-scheme} is a flat formal scheme locally
topologically of finite presentation over $K^\circ$. An \emph{admissible
formal model} of a compact strictly $K$-analytic space $X$ is such a formal
scheme $\mathfrak X$ together with an identification of its generic fibre
$\mathfrak X_\eta$ with $X$.
\end{definition}

The generic-fibre functor in the second definition is Raynaud's. Its
construction and the equivalence after inverting admissible formal blowups
are deep foundations of formal analytic geometry; we use
\cite[Chapter~1]{Berk12}. Algebraic models give formal ones by completion
along the special fibre. Conversely, not every formal model is algebraizable
without extra hypotheses.

\begin{definition}\label{def:model-generic-special-fibres}
For an algebraic model $\mathcal X$, its \emph{generic fibre} and
\emph{special fibre} are
\[
 \mathcal X_K=\mathcal X\times_{K^\circ}\operatorname{Spec}K,\qquad
 \mathcal X_s=\mathcal X\times_{K^\circ}\operatorname{Spec}\widetilde K.
\]
\end{definition}

\begin{example}[The nodal model]
\label{ex:nodal-model}
For a discrete valuation ring $R=K^\circ$, the affine scheme
\[
            \mathcal U=\operatorname{Spec}R[x,y]/(xy-\pi)
\]
is flat because its coordinate ring is torsion-free over $R$. Its generic
fibre is smooth, while its special fibre is the union of the two affine lines
$x=0$ and $y=0$, meeting transversely at the origin. Completion along the
special fibre gives the formal nodal model used in
\cref{chap:four-entrances}.

This is the local equation \eqref{eq:running-node} of the running
semistable degeneration, whose analytic geometry is computed in
\cref{sec:running-semistable}: the generic fibre of $\mathcal U$ contains
the closed annulus $|\pi|\leq|x|\leq1$, and the skeleton of the model is the
interval of Gauss points $v_r$ with $v_r(x)=r\in[0,1]$.  A projective
compactification is the conic degeneration
$\mathcal C=\{XY=\pi Z^2\}\subset\mathbf P^2_R$ of \eqref{eq:running-conic}.
Throughout this chapter we return to $\mathcal U$, to $\mathcal C$, and to
their higher-multiplicity variants $xy=\pi^n$ whenever a new notion needs a
computable test case.
\end{example}

\section{Vertical divisors and divisorial points}

\begin{definition}\label{def:vertical-horizontal-divisors}
A Cartier or Weil divisor on $\mathcal X$ is \emph{vertical} if its support
is contained in $\mathcal X_s$. A divisor is \emph{horizontal} if every
irreducible component dominates $\operatorname{Spec}K^\circ$.
\end{definition}

On a normal noetherian model, every Weil divisor has a unique decomposition
into horizontal and vertical parts. Cartier and Weil divisors must not be
conflated: a vertical prime component need not be Cartier on a singular
model.

The decomposition deserves a complete proof, because it silently uses the
structure of $\operatorname{Spec}K^\circ$ and fails in higher rank.  Recall
that the valuation of $K$ is real-valued, so $K^\circ$ has rank at most one.

\begin{proposition}
\label{prop:vertical-horizontal}
\index{divisor!vertical}\index{divisor!horizontal}
Assume $K$ is nontrivially valued, and let $\mathcal X$ be an algebraic
model of the proper $K$-scheme $X$.
\begin{enumerate}
\item The scheme $\operatorname{Spec}K^\circ$ has exactly two points: the
  open generic point $\eta$ and the closed point $s$.
\item Every irreducible closed subset $Z\subset\mathcal X$ either is
  contained in $\mathcal X_s$, or dominates $\operatorname{Spec}K^\circ$ and
  is the closure of the irreducible closed subset
  $Z_K=Z\cap\mathcal X_K$ of $X$.  The two cases are exclusive.
\item Every irreducible component of $\mathcal X$ is horizontal.
\item The assignments $Z\mapsto Z\cap\mathcal X_K$ and
  $W\mapsto\overline W$ are mutually inverse bijections between the
  horizontal prime divisors of $\mathcal X$ and the prime divisors of $X$.
\item Every Weil divisor $D$ on $\mathcal X$ has a unique decomposition
  $D=D_{\mathrm{hor}}+D_{\mathrm{vert}}$ into a horizontal and a vertical
  part.
\end{enumerate}
\end{proposition}
\begin{proof}
(1) Let $\mathfrak p\ne0$ be a prime of $K^\circ$ and let us show
$\mathfrak p=K^{\circ\circ}$.  Pick a nonzero $a\in\mathfrak p$ and any
$b\in K^{\circ\circ}$.  Since the value group is a subgroup of
$\mathbf R$, it is archimedean: there is $m\geq1$ with
$mv(b)\geq v(a)$, so $a\mid b^m$ by
\Cref{lem:valuation-ideals-principal}, whence
$b^m\in\mathfrak p$ and $b\in\mathfrak p$.  Thus the only primes are $0$
and $K^{\circ\circ}$.  For any nonzero $a\in K^{\circ\circ}$, the basic open set
$D(a)$ consists of the primes not containing $a$, which is exactly
$\{\eta\}$; hence $\{\eta\}$ is open.  Explicitly,
$K^\circ[1/a]=K$: for $x\in K^\times$, the archimedean property of
subgroups of $\mathbf R$ provides $n$ with $nv(a)\geq-v(x)$, so
$xa^n\in K^\circ$.

(2) Let $\zeta$ be the generic point of $Z$ and let
$f:\mathcal X\to\operatorname{Spec}K^\circ$ be the structure morphism.
If $f(\zeta)=s$, then
$f(Z)=f(\overline{\{\zeta\}})\subset\overline{\{s\}}=\{s\}$, so
$Z\subset\mathcal X_s$.  If $f(\zeta)=\eta$, then
$Z_K=Z\cap\mathcal X_K$ is a nonempty open subset of $Z$, because
$\mathcal X_K=f^{-1}(\eta)$ is open in $\mathcal X$ by (1).  A nonempty
open subset of an irreducible space is irreducible and dense, so $Z$ is its
closure and $Z$ dominates $\operatorname{Spec}K^\circ$.  The cases exclude
each other because a subset of $\mathcal X_s$ maps to $\{s\}$.

(3) Let $Z$ be an irreducible component with generic point $\zeta$ and
suppose $Z\subset\mathcal X_s$.  Then the image of $\pi'$ in
$\mathcal O_{\mathcal X,\zeta}$ lies in the maximal ideal for every
$\pi'\in K^{\circ\circ}$.  Fix a nonzero $a\in K^{\circ\circ}$.  Since
$\zeta$ is a generic point of $\mathcal X$, the local ring
$\mathcal O_{\mathcal X,\zeta}$ has a unique prime, so its maximal ideal is
its nilradical, and $a$ is nilpotent in it: $a^n=0$ in
$\mathcal O_{\mathcal X,\zeta}$ for some $n$.  Hence there is a section $g$
of $\mathcal O_{\mathcal X}$ on some affine chart with $g\notin\zeta$, in
particular $g\ne0$, and $a^ng=0$.  Then $g$ is a nonzero torsion section,
contradicting the flatness of $\mathcal X$ via
\Cref{lem:torsion-free-flat-valuation}.

(4) Let $Z$ be a horizontal prime divisor with generic point $\zeta$.  The
local rings of $\mathcal X$ and of the open subscheme $\mathcal X_K$ at
$\zeta$ coincide, so $Z_K$ has codimension one in $X$ exactly when $Z$ has
codimension one in $\mathcal X$.  Conversely, if $W\subset X$ is a prime
divisor, its closure $\overline W\subset\mathcal X$ is irreducible,
horizontal, satisfies $\overline W\cap\mathcal X_K=W$ because $W$ is
already closed in $\mathcal X_K$, and has codimension one by the same
equality of local rings.  The two maps are inverse to each other by (2).

(5) Group the prime components of $D$ according to the dichotomy in (2);
by (2) no prime divisor is simultaneously vertical and horizontal, so the
decomposition is unique.
\end{proof}

\begin{remark}[Failure in higher rank]
\label{rem:vertical-horizontal-higher-rank}
\index{valuation ring!higher rank}
For a valuation ring $R$ of rank two the trichotomy is genuinely worse, and
this is one reason the chapter insists on real-valued valuations.  Let
$\Gamma=\mathbf Z\times\mathbf Z$ be ordered lexicographically, and let $R$
be a valuation ring with value group $\Gamma$, for instance the valuation
ring of $v$ on $F=k(s,t)$ determined by $v(s)=(1,0)$ and $v(t)=(0,1)$.
Then $\operatorname{Spec}R=\{0\subset\mathfrak p\subset\mathfrak m\}$ has
three points, where $\mathfrak p$ corresponds to the convex subgroup
$0\times\mathbf Z$.  In $\mathbf A^1_R=\operatorname{Spec}R[T]$, the closed
subset $V(\mathfrak pR[T])=\operatorname{Spec}(R/\mathfrak p)[T]$ is
irreducible of codimension one.  It is not vertical, since it contains
points lying over $\mathfrak p$, and it does not dominate
$\operatorname{Spec}R$, since its image is
$\{\mathfrak p,\mathfrak m\}$.  A prime divisor over such a base can
therefore be neither horizontal nor vertical, and the decomposition in
\Cref{prop:vertical-horizontal} must be replaced by a stratification along
$\operatorname{Spec}R$.  None of the constructions of this book requires
that generality.
\end{remark}

Assume for the remainder of this section that $K$ is discretely valued,
$X$ is integral, and $\mathcal X$ is normal. Write the Cartier divisor
of $\pi$ as
\[
                       \operatorname{div}_{\mathcal X}(\pi)
                       =\mathcal X_s=\sum_E b_EE.
\]
Here $E$ runs over the irreducible components of the special fibre and
$b_E=\operatorname{ord}_E(\pi)>0$.

\begin{proposition}
\label{prop:component-valuation}
For each component $E$, the formula
\[
                         v_E(f)=\frac{1}{b_E}\operatorname{ord}_E(f)
\]
defines a valuation on $K(X)$ extending $v_K$. Its center on $\mathcal X$ is
the generic point of $E$.
\end{proposition}
\begin{proof}
Normality implies that the local ring at the generic point of $E$ is a
discrete valuation ring, so $\operatorname{ord}_E$ is a valuation. Since
$\operatorname{ord}_E(\pi)=b_E$, the normalized valuation takes $\pi$ to $1$
and hence restricts to $v_K$ on $K$. Its valuation ring dominates the local
ring at the generic point of $E$, which is exactly the assertion about its
center.
\end{proof}

The analytic point associated to $v_E$ is denoted $x_E$ and is called a
\emph{divisorial point}. If $D=\sum_E a_EE$ is a vertical Cartier divisor,
then its local equation $f_D$ defines a model function
\[
                         \varphi_D(x)=-\log|f_D(x)|=v_x(f_D).
\]
At a divisorial point,
\[
                         \varphi_D(x_E)=\frac{a_E}{b_E}.
\]
This formula fixes our sign convention: an effective vertical divisor gives
a nonnegative model function, and $\varphi_{\mathcal X_s}=1$.

\begin{lemma}
\label{lem:model-function-well-defined}
The local formula $\varphi_D=-\log|f_D|$ glues to a continuous function on the
generic Berkovich space, and pullback of $D$ to a dominating model defines the
same function.
\end{lemma}
\begin{proof}
On an overlap, two local equations differ by a unit $u$ in the model.
Both $u$ and $u^{-1}$ are power-bounded on the generic fibre, hence
$|u(x)|=1$ there. Thus $-\log|f_D|$ is independent of the local equation.
After pullback, a local equation becomes its pullback, whose value at an
analytic point is unchanged by functoriality of evaluation. Continuity follows
on each formal affine chart because the absolute value of an analytic
function is continuous.
\end{proof}

\begin{example}[Divisorial points and model functions on the nodal model]
\label{ex:node-divisorial-model-function}
\index{divisorial point!nodal model}\index{model function!example}
Work on the nodal model
$\mathcal U=\operatorname{Spec}R[x,y]/(xy-\pi)$ of \Cref{ex:nodal-model},
whose total space is regular.  The special fibre is
$E_x\cup E_y$ with $E_x=\{x=0\}$ and $E_y=\{y=0\}$.

At the generic point of $E_x$ the coordinate $y$ is invertible, so
$\pi=xy$ has order $\operatorname{ord}_{E_x}(\pi)=\operatorname{ord}_{E_x}(x)=1$;
hence $b_{E_x}=1$, and symmetrically $b_{E_y}=1$.  The divisorial point
$x_{E_x}$ is the valuation with
\[
                 v_{E_x}(x)=1,\qquad v_{E_x}(y)=0,
\]
which is the Gauss point $v_1$ of the skeleton computed in
\cref{chap:four-entrances} and \cref{sec:running-semistable}; likewise
$x_{E_y}=v_0$.  The two divisorial points of the model are exactly the two
endpoints of the skeleton, and the parameter $r=v(x)$ interpolates between
them through monomial valuations.

Now take the vertical divisor $D=E_x$.  On the regular scheme $\mathcal U$
it is Cartier: since $x$ vanishes precisely on $E_x$
(the equation $x=0$ forces $\pi=xy=0$), with multiplicity one, a global
equation is $f_D=x$.  The associated model function is
\[
     \varphi_{E_x}(v_r)=v_r(x)=r,
\]
the affine function of slope one on the skeleton $[0,1]$, vanishing at
$v_0$ and equal to $1=a_{E_x}/b_{E_x}$ at $v_1$.  This confirms both sign
tests fixed in the Conventions: the effective divisor $E_x$ produces a
nonnegative function, and
$\varphi_{\mathcal U_s}=\varphi_{E_x}+\varphi_{E_y}=v_r(x)+v_r(y)=1$.
Restrictions of model functions to skeleta are studied systematically in
\cref{sec:skeleton-model-functions}.
\end{example}

\section{Domination and common refinements}

\begin{definition}\label{def:model-domination}
A model $\mathcal Y$ \emph{dominates} a model $\mathcal X$ if there is a
proper $K^\circ$-morphism
$\rho:\mathcal Y\to\mathcal X$ whose generic fibre is the identity of $X$.
Such a morphism is also called a \emph{vertical modification}. Models with
domination form a preorder; after identifying isomorphic models they form a
partially ordered system.
\end{definition}

\paragraph{Four words that are not synonyms.}
\index{modification}\index{domination}\index{refinement}\index{model!versus modification}
The Conventions promise that \emph{modification}, \emph{model},
\emph{domination}, and \emph{refinement} are never used interchangeably.
This is the natural place to fix the distinctions.
\begin{enumerate}
\item A \emph{modification} is a proper birational morphism of integral
  schemes (\Cref{def:appB-modification}).  It carries no reference to a
  base; both source and target live in a single birational class.
\item A \emph{model} is an object over $K^\circ$ \emph{together with} a
  marked isomorphism of its generic fibre with $X$.  A model is not a
  morphism, and an abstract $K^\circ$-scheme with the right generic fibre
  is not yet a model until the identification is chosen.
\item \emph{Domination} is the morphism notion for models: a proper
  $K^\circ$-morphism restricting to the identity on $X$.  When $X$ is
  integral, the generic fibres are dense in both flat models by
  \Cref{prop:vertical-horizontal}, so a domination is automatically a
  modification of total spaces over $K^\circ$; the converse fails, since a
  modification need not respect the marked identifications.
\item A \emph{refinement} is a combinatorial notion for skeleta and dual
  complexes: a subdivision of the underlying polyhedral structure
  (\cref{sec:skeleton-refinement}).  A domination of SNC models induces a
  refinement-type comparison of skeleta, but domination is a geometric
  relation and refinement a combinatorial one, and neither determines the
  other.
\end{enumerate}
The trivially valued theory adds a fifth word, \emph{test configuration},
for a degeneration over $\mathbf A^1$ with a $\mathbf G_m$-action; it is
related to, but not a special case of, the models of this chapter, and is
treated in \cref{chap:bj-test-configurations}.  The situation of a common
domination is summarized by the commutative diagram
\[
\begin{tikzcd}[column sep=large,row sep=large]
 & \mathcal Y \arrow[dl,"\rho_1"'] \arrow[dr,"\rho_2"] & \\
\mathcal X_1 & X \arrow[l,hook'] \arrow[u,hook] \arrow[r,hook]
 & \mathcal X_2
\end{tikzcd}
\]
in which every hooked arrow is the inclusion of the generic fibre through
the marked identification, and both triangles restrict to the identity of
$X$ over the generic point.  The next results construct $\mathcal Y$.

We first construct the schematic closure of a generic-fibre subscheme in
complete generality.  No noetherian hypothesis is needed; the only standing
assumption is that $K$ is nontrivially valued, so that the generic point of
$\operatorname{Spec}K^\circ$ is open
(\Cref{prop:vertical-horizontal}\,(1)).

\begin{lemma}
\label{lem:flat-closure-generic-fibre}
\index{flat closure}\index{schematic closure}
Let $\mathcal P$ be a $K^\circ$-scheme and let $W\subset\mathcal P_K$ be a
closed subscheme of its generic fibre.  There is a unique closed subscheme
$\overline W\subset\mathcal P$ that is flat over $K^\circ$ and satisfies
$\overline W\cap\mathcal P_K=W$ as subschemes of $\mathcal P_K$.  Its
underlying set is the topological closure of $W$, and every closed
subscheme of $\mathcal P$ whose generic fibre contains $W$ contains
$\overline W$.  We call $\overline W$ the \emph{flat closure} of $W$.
\end{lemma}
\begin{proof}
Fix a nonzero $a\in K^{\circ\circ}$, so that $K=K^\circ[1/a]$ and, for
every affine chart $\operatorname{Spec}A\subset\mathcal P$, the
corresponding chart of the generic fibre is
$\operatorname{Spec}A_K$ with $A_K=A\otimes_{K^\circ}K=A[1/a]$.

\emph{Construction.}  On a chart $\operatorname{Spec}A$, let
$I_W\subset A_K$ be the ideal of $W$ and let
\[
             I=\ker\bigl(A\longrightarrow A_K/I_W\bigr),
\]
the contraction of $I_W$.  Localization at an element $g\in A$ is exact and
commutes with the passage from $A$ to $A_K$, so the formation of $I$
commutes with localization: $I_g=\ker(A_g\to (A_g)_K/(I_W)_g)$.  Any two
affine charts of $\mathcal P$ are covered by opens that are simultaneously
distinguished in both, so the chartwise ideals agree on overlaps and define
a quasi-coherent ideal sheaf.  Let $\overline W\subset\mathcal P$ be the
closed subscheme it cuts out.

\emph{Flatness and generic fibre.}  On each chart, $A/I$ injects into the
$K$-vector space $A_K/I_W$ and is therefore torsion-free, hence flat by
\Cref{lem:torsion-free-flat-valuation}.  For the generic fibre we must
check that $I\cdot A_K=I_W$.  The inclusion $\subset$ is clear.  For the
converse, let $f\in I_W$.  Since $A_K=A[1/a]$, there is an $n$ with
$a^nf\in A$; the element $a^nf$ maps into $I_W$, hence lies in $I$, and
$f=a^{-n}(a^nf)\in I\cdot A_K$.  Thus
$(A/I)\otimes_{K^\circ}K=A_K/I_W$, that is,
$\overline W\cap\mathcal P_K=W$.

\emph{Minimality and uniqueness.}  Let $\mathcal Z=V(J)$ be a closed
subscheme with $W\subset\mathcal Z\cap\mathcal P_K$, i.e.\
$J\cdot A_K\subset I_W$ on each chart.  Then $J$ maps to zero in
$A_K/I_W$, so $J\subset I$ and $\overline W\subset\mathcal Z$.  If
moreover $\mathcal Z$ is flat with $\mathcal Z\cap\mathcal P_K=W$, then
$A/J$ is torsion-free, so $A/J\to(A/J)\otimes_{K^\circ}K=A_K/I_W$ is
injective, which says exactly $J=I$.

\emph{Support.}  The set $V(I)\cap D(a)$ equals $V(I_W)=W$, and it is
dense in $V(I)$: if $\mathfrak p$ is a minimal prime of $A/I$ containing
$a$, then $a$ is nilpotent in the one-prime local ring $(A/I)_{\mathfrak
p}$, so $sa^n=0$ in $A/I$ for some $s\notin\mathfrak p$; the element $s$ is
then a nonzero torsion element of the torsion-free module $A/I$, a
contradiction.  Hence no component of $\overline W$ lies over the closed
point, and $|\overline W|$ is the closure of $|W|$.
\end{proof}

\begin{proposition}\label{prop:common-domination}
Any two proper algebraic models $\mathcal X_1,\mathcal X_2$ of $X$ admit a
common algebraic model $\mathcal Y$ dominating both.  It may be chosen as
the flat closure of the diagonal generic fibre inside
$\mathcal X_1\times_{K^\circ}\mathcal X_2$.  If the given models are
projective, then $\mathcal Y$ may be chosen projective.  If $X$ is integral
and normal and normalization is finite in the category under consideration,
then $\mathcal Y$ may also be chosen normal.
\end{proposition}
\begin{proof}
If the valuation is trivial, then $K^\circ=K$, every model is its own generic
fibre, and the marked identifications make $X$ itself a common model.  We may
therefore assume that the valuation is nontrivial.  We spell out the
graph-closure construction in five steps.  Write
$j_i:X\xrightarrow{\ \sim\ }(\mathcal X_i)_K\subset\mathcal X_i$ for the
two marked identifications followed by the open immersions of the generic
fibres (\Cref{prop:vertical-horizontal}\,(1)), and set
$\mathcal P=\mathcal X_1\times_{K^\circ}\mathcal X_2$.

\emph{Step 1: the graph embedding.}  The morphism
$\gamma=(j_1,j_2):X\to\mathcal P$ factors as
\[
   X\xrightarrow{\ \Gamma\ }
   X\times_K(\mathcal X_2)_K
   =\mathcal P_K\cap\bigl((\mathcal X_1)_K\times_{K^\circ}\mathcal X_2\bigr)
   \lhook\joinrel\longrightarrow\mathcal P,
\]
where $\Gamma$ is the graph of $j_2$ viewed over $K$, under the
identification $(\mathcal X_1)_K\simeq X$.  Since $\mathcal X_2$ is proper,
hence separated, over $K^\circ$, the graph $\Gamma$ is a closed immersion,
and $\mathcal P_K=(\mathcal X_1)_K\times_K(\mathcal X_2)_K$ is open in
$\mathcal P$.  Thus $\gamma$ identifies $X$ with a closed subscheme
$W\subset\mathcal P_K$.

\emph{Step 2: closure.}  Let $\mathcal Y=\overline W$ be the flat closure
of \Cref{lem:flat-closure-generic-fibre}.  It is a closed subscheme of
$\mathcal P$, flat over $K^\circ$, with generic fibre $W\simeq X$.
Properness of $\mathcal X_1$ and $\mathcal X_2$ makes $\mathcal P$ proper
over $K^\circ$, and a closed subscheme of a proper scheme is proper.

\emph{Step 3: the projections dominate.}  Let
$\rho_i=\operatorname{pr}_i|_{\mathcal Y}:\mathcal Y\to\mathcal X_i$.  Each
$\rho_i$ is proper: the product projection
$\mathcal P\to\mathcal X_i$ is the base change of the proper morphism
$\mathcal X_j\to\operatorname{Spec}K^\circ$ for $j\ne i$, and
$\mathcal Y\hookrightarrow\mathcal P$ is a closed immersion.  On the
generic fibre, $\rho_i$ restricts to
$\operatorname{pr}_i\circ\gamma=j_i$, which is the marked identification;
so $\rho_i$ is a domination of models.

\emph{Step 4: the projective case.}  If both $\mathcal X_i$ are projective
over $K^\circ$, so is $\mathcal P$ by the Segre embedding, and the closed
subscheme $\mathcal Y$ is projective.  This proves the third assertion.

\emph{Step 5: finiteness and normality.}  If $K^\circ$ is noetherian, i.e.\
a discrete valuation ring or a field
(\Cref{prop:noetherian-valuation-dvr}), then $\mathcal P$ is of finite
type over $K^\circ$, hence noetherian, and the closed subscheme $\mathcal
Y$ is of finite type; over a noetherian base, finite type and finite
presentation agree, so $\mathcal Y$ is an algebraic model.  Without the
noetherian hypothesis, finite presentation of this same flat closure is
proved in \Cref{prop:flat-closure-finite-presentation} below.  Finally,
suppose $X$ is integral and normal, and that normalization is finite in
the category at hand---for instance, $K^\circ$ a complete discrete
valuation ring, which is excellent \cite[Scholie~7.8.3]{EGAIV2}, so that
finite-type $K^\circ$-schemes are Nagata and
\cite[Tag 035S]{stacks-project} applies.  Since $X$ is irreducible and
reduced, its closure $\mathcal Y$ is integral: irreducibility passes to
closures, and on charts $A/I$ embeds into the reduced ring $A_K/I_W$ of
functions on $X$.  The normalization
$\nu:\mathcal Y^\nu\to\mathcal Y$ is then finite, so $\mathcal Y^\nu$ is
proper over $K^\circ$.  On every affine chart its coordinate ring injects
into its function field, in which every nonzero element of $K^\circ$ remains
nonzero; it is therefore $K^\circ$-torsion-free and hence flat by
\Cref{lem:torsion-free-flat-valuation}.  Over the generic fibre, $\nu$ is an
isomorphism: the local rings of $\mathcal Y$ at points of
$\mathcal Y_K\simeq X$ are local rings of the normal scheme $X$, and
normalization does not change an integral scheme over its normal locus.
Hence $\mathcal Y^\nu$ is again a model of $X$, normal, and dominates both
$\mathcal X_i$ through $\nu$.
\end{proof}

\begin{remark}\label{rem:normal-common-domination-scope}
The normality assertion requires $X$ itself to be normal: normalizing a
model changes the generic fibre into the normalization of $X$, which is a
model of $X$ only when $X$ was already normal.  For non-normal $X$, the
correct statement is that $\mathcal Y^\nu$ is a normal model of the
normalization $X^\nu$; the interplay between the two levels is exactly the
subject of the descent chapters of \cref{part:bj}.
\end{remark}

\begin{remark}\label{rem:common-domination-finite-presentation}
The graph-closure argument itself first produces a proper flat scheme.  Over
an arbitrary nonnoetherian valuation ring, finite presentation is a separate
conclusion rather than a formal property of a closed subscheme.  The next
proposition supplies precisely this missing step, so the object asserted in
\Cref{prop:common-domination} is an algebraic model over every valuation
ring covered by the chapter.
\end{remark}

The required approximation theorem is in fact available with the tools of
this chapter: it is an application of flattening by blowup, whose
Stacks Project form \cite[Tag 0815]{stacks-project} assumes only a
quasi-compact quasi-separated base and therefore applies to
$\operatorname{Spec}K^\circ$ for every valuation ring.  Over a valuation
ring the blowup produced by the flattening theorem is trivial, and only its
finite-presentation conclusion survives---which is exactly what is missing.

\begin{proposition}
\label{prop:flat-closure-finite-presentation}
\index{model!finite presentation}\index{flattening!over a valuation ring}
Assume $K$ is nontrivially valued.  Let $\mathcal P$ be a finitely
presented $K^\circ$-scheme and let $W\subset\mathcal P_K$ be a closed
subscheme of its generic fibre.  Then the flat closure
$\overline W\subset\mathcal P$ of \Cref{lem:flat-closure-generic-fibre} is
finitely presented over $K^\circ$.  In particular, in
\Cref{prop:common-domination} the common dominating model $\mathcal Y$ is
always an algebraic model, with no noetherian hypothesis on $K^\circ$.
\end{proposition}
\begin{proof}
\emph{Step 1: the flattening input.}  Apply
\cite[Tag 0815]{stacks-project} with the following data: the base
$S=\operatorname{Spec}K^\circ$, which is affine, hence quasi-compact and
quasi-separated; the quasi-compact open $U=\{\eta\}=D(a)$ for a fixed
nonzero $a\in K^{\circ\circ}$ (\Cref{prop:vertical-horizontal}\,(1)); the
quasi-compact, locally finitely presented $S$-scheme $\mathcal P$; and the
finite-type quasi-coherent module
$\mathcal F=\mathcal O_{\overline W}=\mathcal O_{\mathcal P}/\mathcal I$,
a cyclic quotient of the structure sheaf.  Over $U$ the restriction of
$\mathcal F$ is $\mathcal O_W$, which is a coherent, hence finitely
presented, module on the noetherian scheme $\mathcal P_K$, and it is flat
over the field $K$.  The theorem produces a $U$-admissible blowup
$S'\to S$ such that the strict transform of $\mathcal F$ is finitely
presented and flat over $S'$.

\emph{Step 2: the blowup is trivial.}  The center of the blowup is a
finite-type quasi-coherent ideal of $K^\circ$ that restricts to the unit
ideal on $U$; by \Cref{lem:valuation-ideals-principal} it is principal,
say $(b)$, and $b\ne0$ because the ideal is nonzero on $U$.  The Rees
algebra $\bigoplus_{m\geq0}(b)^m$ is the polynomial algebra on the
degree-one generator $b$, since $b$ is a nonzerodivisor; hence
$S'=\operatorname{Proj}\bigoplus_m(b)^m\to S$ is an isomorphism.

\emph{Step 3: the strict transform is $\mathcal F$ itself.}  Under the
identity blowup, the strict transform of $\mathcal F$ is the quotient of
$\mathcal F$ by its subsheaf of sections supported on $V(b)$.  A local
section $s$ of $\mathcal F$ supported on $V(b)$ has zero stalk at every
point of the generic fibre over the chart in question, and since
$\mathcal F$ is of finite type this means $b^ns=0$ locally for some $n$;
in particular $s$ is a $K^\circ$-torsion section.  But
$\mathcal O_{\overline W}$ is flat, hence torsion-free, so $s=0$.  The
strict transform is therefore $\mathcal F$ itself, and Step 1 says that
$\mathcal F=\mathcal O_{\mathcal P}/\mathcal I$ is finitely presented as an
$\mathcal O_{\mathcal P}$-module.

\emph{Step 4: from module presentation to scheme presentation.}  We claim
$\mathcal I$ is of finite type.  Locally on a chart
$\operatorname{Spec}A$, we have a surjection $q:A\to M=A/I$ with $M$
finitely presented, say with presentation
$A^r\xrightarrow{\varphi}A^s\xrightarrow{\psi}M\to0$.  Schanuel's lemma
applied to the two exact sequences
$0\to I\to A\to M\to0$ and
$0\to\ker\psi\to A^s\to M\to0$
gives $I\oplus A^s\simeq\ker\psi\oplus A$, and $\ker\psi=\operatorname{im}\varphi$
is finitely generated; hence $I$, a quotient of
$\ker\psi\oplus A$, is finitely generated.  Thus
$\overline W\hookrightarrow\mathcal P$ is a closed immersion of finite
presentation, and the composition
$\overline W\to\mathcal P\to\operatorname{Spec}K^\circ$ is finitely
presented.

For the final assertion, take $\mathcal P=\mathcal X_1\times_{K^\circ}
\mathcal X_2$, which is finitely presented because both factors are; then
$\mathcal Y=\overline W$ is proper, flat, and finitely presented, i.e.\ an
algebraic model.
\end{proof}

\begin{remark}\label{rem:flat-closure-raynaud-gruson}
\Cref{prop:flat-closure-finite-presentation} is a special case of the
theorem of Raynaud--Gruson that a flat finite-type module over an integral
scheme is automatically of finite presentation; both statements originate
in \cite{RaynaudGruson1971}.  We derived it from
\cite[Tag 0815]{stacks-project} to keep the book's external dependencies
within the flattening theorem already recorded in \cref{app:stacks},
and because the derivation shows exactly where the valuation-ring
structure enters: through the principality of finitely generated ideals.
Note also that the noetherian hypotheses in the model-flattening
\Cref{thm:flattening-models} below belong to that statement's stronger
strict-transform conclusion, not to the tag itself.
\end{remark}

\begin{example}[Models over a densely valued field]
\label{ex:model-dense-value-group}
\index{model!nonnoetherian base}\index{value group!dense}
Let $K$ be a complete algebraically closed nontrivially valued field, for
example the completion of an algebraic closure of $\mathbf Q_p$.  Its value
group $\Gamma=v(K^\times)$ is divisible, hence dense in $\mathbf R$, and
$K^\circ$ is not noetherian by \Cref{prop:noetherian-valuation-dvr}.
Concretely, choose $a_n\in K^{\circ\circ}$ with $v(a_n)$ strictly
decreasing to $0$; then $(a_1)\subsetneq(a_2)\subsetneq\cdots$ is a
nonstabilizing chain of ideals, and
$K^{\circ\circ}=\bigcup_n(a_n)$ is not finitely generated.

Models nevertheless exist and behave well.  Take
$X=\mathbf P^1_K$ and $\mathcal X=\mathbf P^1_{K^\circ}$, a flat, proper,
finitely presented model.  Fix $a\in K^{\circ\circ}$ nonzero and blow up
the finitely generated vertical ideal $(T,a)$ on the affine chart
$\operatorname{Spec}K^\circ[T]$, glued trivially elsewhere.  The two
standard charts are
\[
 \operatorname{Spec}K^\circ[T,s]/(a-Ts),
 \qquad
 \operatorname{Spec}K^\circ[u]\quad(T=au),
\]
and the resulting model $\mathcal X_a$ is again flat, proper, and finitely
presented; its special fibre acquires a node with local equation $Ts=a$,
and the skeleton of $\mathcal X_a$ is a segment of length $v(a)$, exactly
as for the nodal model over a discrete valuation ring.

What fails is every statement that mentions a uniformizer or noetherian
finiteness:
\begin{enumerate}
\item there is no element $\pi$ of minimal positive value, so
  $\mathcal X_s$ is not the divisor of a distinguished function; each
  nonzero $a\in K^{\circ\circ}$ gives its own vertical divisor
  $\operatorname{div}(a)$, with $\varphi_{\operatorname{div}(a)}=v(a)$
  constant, and only the normalized classes
  $v(a)^{-1}\operatorname{div}(a)$ are comparable;
\item multiplicities $b_E=\operatorname{ord}_E(\pi)$ are meaningless as
  stated, and the generic local ring of a component need not be a discrete
  valuation ring.  Components of an admissible formal model are instead
  read through their reduction tubes and Shilov boundary points.  In the
  discrete normal case this recovers
  $b_E^{-1}\operatorname{ord}_E$; no such formula is used in the dense-value
  case;
\item the total space $\mathcal X$ is not noetherian: already
  $K^\circ[T]$ contains the nonstabilizing chain
  $(a_1)\subsetneq(a_2)\subsetneq\cdots$, so the noetherian intersection
  theory of this chapter does not literally apply, and Chambert--Loir
  masses over such fields are constructed by the limit and base-change
  arguments of \cref{chap:metrics};
\item only coherent finitely generated centers remain in the finitely
  presented model category: the ideal
  $K^{\circ\circ}\mathcal O_{\mathcal X}$, for example, is not coherent and
  hence is not admissible,
  and the cofiltered system of models must be indexed by finitely
  presented data.
\end{enumerate}
On the positive side, \Cref{lem:torsion-free-flat-valuation},
\Cref{lem:flat-closure-generic-fibre}, and
\Cref{prop:flat-closure-finite-presentation} are insensitive to
noetherianity: common dominations of finitely presented models exist over
every such $K^\circ$.  This division of labour---flatness and closure
arguments valid always, multiplicity and intersection arguments valid over
a discrete valuation ring---is the reason the chapter states its standing
hypotheses section by section.
\end{example}

\begin{definition}\label{def:vertical-admissible-blowup}
Let $\mathcal I\subset\mathcal O_{\mathcal X}$ be a coherent ideal whose
restriction to the generic fibre is the unit ideal. The blowup
$\operatorname{Bl}_{\mathcal I}\mathcal X\to\mathcal X$ is a
\emph{vertical blowup}. In the formal category, a coherent open ideal is
\emph{admissible} if locally it contains an ideal of definition; its formal
blowup is an \emph{admissible formal blowup}.
\end{definition}

\begin{proposition}
\label{prop:vertical-blowup-generic-iso}
A vertical blowup is an isomorphism on the generic fibre. If $\mathcal X$ is
projective, its blowup is projective over $K^\circ$.
\end{proposition}
\begin{proof}
Blowup commutes with flat base change. Since $\mathcal I_K=\mathcal O_X$, the
generic fibre is the blowup of the unit ideal, hence is $X$. The blowup is
$\operatorname{Proj}_{\mathcal X}\bigoplus_{m\geq0}\mathcal I^m$, so it is
projective over $\mathcal X$; composition with a projective structure
morphism remains projective.
\end{proof}

\begin{example}[Blowing up a vertical Weil divisor]
\label{ex:vertical-blowup-weil-center}
\index{blowup!vertical}\index{divisor!Weil versus Cartier}
Let $R=K^\circ$ be a discrete valuation ring and let
\[
        \mathcal X=\operatorname{Spec}A,\qquad A=R[x,y]/(xy-\pi^2).
\]
The total space has a single singular point, the origin of the special
fibre, of type $A_1$.  The special fibre is $D_x\cup D_y$ with
$D_x=\{x=\pi=0\}$ and $D_y=\{y=\pi=0\}$, both of multiplicity one.  The
prime divisor $D_x$ is \emph{not} Cartier at the singular point: in the
local ring there, the ideal $(x,\pi)$ of $D_x$ needs both generators.
(Compare \Cref{ex:appF-non-cartier-bdivisor} for the same phenomenon in a
purely birational setting.)  So a vertical prime component of a singular
model is a Weil divisor with no local equation, and blowing it up is a
genuine operation.

Blow up the ideal $\mathcal I=(x,\pi)$.  On the chart where $x$ generates,
introduce $s=\pi/x$: the relations $\pi=xs$ and $xy=\pi^2=x^2s^2$ give
$y=xs^2$, and the chart is
\[
        \operatorname{Spec}R[x,s]/(\pi-xs),
\]
a regular scheme with semistable special fibre $\{xs=0\}$.  On the chart
where $\pi$ generates, introduce $u=x/\pi$: from $xy=\pi^2$ we get
$uy=\pi$, and the chart is
\[
        \operatorname{Spec}R[u,y]/(uy-\pi),
\]
again regular and semistable.  The gluing is $u=1/s$.  The special fibre of
the blowup is a chain
\[
             \widetilde D_y\ \cup\ E\ \cup\ \widetilde D_x ,
\]
where $E$, the component with charts $\{x=0\}$ and $\{y=0\}$ above the
center, is a projective line: its two affine pieces
$\operatorname{Spec}\widetilde K[s]$ and
$\operatorname{Spec}\widetilde K[u]$ glue along $u=1/s$.  All
multiplicities are one, and the two nodes have local equations $xs=\pi$
and $uy=\pi$.  Thus a single vertical blowup with a one-dimensional
non-Cartier center resolves the $A_1$-model into a strictly semistable
SNC model; the skeleton grows from the segment $[0,2]$ with endpoints only
to the same segment subdivided at the midpoint $v(x)=1$, the divisorial
point of $E$.
\end{example}

Normalized blowups are obtained by normalizing the blowup.
\index{blowup!normalized}
Over an excellent
discrete valuation ring and for finite type models, normalization is finite:
a complete discrete valuation ring is excellent
\cite[Scholie~7.8.3]{EGAIV2}, excellence passes to finite-type algebras, and
the normalization of a Nagata scheme is finite by
\cite[Tag 035S]{stacks-project}.
In greater generality this finiteness is a hypothesis, not a convention.
When the generic fibre $X$ is normal, the composite of blowup and
normalization is again a vertical modification,
projective when the model is
(\Cref{prop:vertical-blowup-generic-iso} together with finiteness of
$\nu$); the systematic theory, including the Rees valuations that the
exceptional components of a normalized blowup realize, is developed in
\cref{sec:appB-normalized-blowups}.
Raynaud's theorem says that admissible formal blowups become isomorphisms on
generic fibres and that, after inverting them, admissible formal models encode
quasi-compact admissible analytic spaces. We use this theorem in
\Cref{chap:berkovich} rather than reproving formal GAGA.

\begin{example}[A normalized blowup computed on a two-dimensional model]
\label{ex:normalized-blowup-computed}
\index{blowup!normalized!example}
Let $R=K^\circ$ be a discrete valuation ring, take the regular model
$\mathcal X=\mathbf A^1_R=\operatorname{Spec}R[T]$ of $\mathbf A^1_K$, and
blow up the vertical ideal $\mathfrak a=(T^2,\pi^2)$.  On the chart of
$\operatorname{Bl}_{\mathfrak a}\mathcal X$ where $T^2$ generates, set
$w=\pi^2/T^2$; the chart is
\[
    \operatorname{Spec}R[T,w]/(\pi^2-T^2w).
\]
This scheme is \emph{not normal}.  At the height-one prime
$\mathfrak p=(T,\pi)$, the maximal ideal of the one-dimensional local ring
is generated by the classes of $T$ and $\pi$, while the relation
$\pi^2-T^2w$ has no linear term.  Thus
$\dim_{\kappa(\mathfrak p)}\mathfrak p/\mathfrak p^2=2$, so the local ring
is not regular and hence is not a discrete valuation ring.  A
one-dimensional noetherian normal local domain is a discrete valuation
ring, proving nonnormality without a residue-characteristic assumption.

Normalize.  The rational function $t=\pi/T$ satisfies
$t^2=\pi^2/T^2=w$, an integral equation over the chart ring, so $t$ lies
in the normalization.  Adjoining it and eliminating $w=t^2$ yields
\[
      \operatorname{Spec}R[T,t]/(\pi-Tt),
\]
which is regular, hence normal, and finite over the chart (it is generated
by the single integral element $t$).  Therefore the normalized blowup of
$(T^2,\pi^2)$ is the standard semistable model with node $Tt=\pi$: the
same model as the blowup of the ideal $(T,\pi)$.  This is no accident: the
integral closure of $(T^2,\pi^2)$ is $(T,\pi)^2=(T^2,T\pi,\pi^2)$, and a
normalized blowup depends only on the integral closure of its center
(\cref{sec:appB-normalized-blowups}).  The exceptional component $E$ has
$\operatorname{ord}_E(T)=\operatorname{ord}_E(\pi)=1$, so its divisorial
point is the Gauss point of the disc of radius $|\pi|$; it realizes the
unique Rees valuation of $\mathfrak a$ in the sense of
\Cref{thm:appB-rees-valuations}.
\end{example}

\begin{example}[A non-normal model and its normalization]
\label{ex:non-normal-model-normalization}
\index{model!non-normal}\index{normalization!of a model}
Non-normal models also arise directly, not only as intermediate blowups.
Let $R=K^\circ$ be a discrete valuation ring and
\[
        \mathcal Z=\operatorname{Spec}B,\qquad B=R[x,y]/(x^2y-\pi^2).
\]
The polynomial $x^2y-\pi^2$ is irreducible in the unique factorization
domain $R[x,y]$ (it is primitive of degree one in $y$), so $B$ is a domain,
torsion-free, and $\mathcal Z$ is a flat affine model.  Its generic fibre
is smooth: in $B\otimes_RK$ the element $x$ is invertible, since
$x\cdot(xy/\pi^2)=1$, and
$B\otimes_RK\simeq K[x,x^{-1}]$ is the coordinate ring of
$\mathbf G_{m,K}$.  Its special fibre, however, is
$\operatorname{Spec}\widetilde K[x,y]/(x^2y)$: a \emph{doubled} line
$\{x^2=0\}$ together with the reduced line $\{y=0\}$.

At the height-one prime $\mathfrak p=(x,\pi)$, the maximal ideal of the
one-dimensional local ring $B_{\mathfrak p}$ is generated by $x$ and
$\pi$, while the relation $\pi^2-yx^2$ lies in
$\mathfrak p^2$.  Hence
$\dim_{\kappa(\mathfrak p)}\mathfrak p/\mathfrak p^2=2$, so
$B_{\mathfrak p}$ is not a discrete valuation ring and $\mathcal Z$ is not
normal.  Note the practical consequence: at the generic point of
$\{x=\pi=0\}$, $\operatorname{ord}_E$ is not defined and
\Cref{prop:component-valuation} genuinely needs its normality hypothesis.

The normalization is computed by hand.  The function $t=\pi/x$ satisfies
$t^2=\pi^2/x^2=y$, so $t$ is integral over $B$, and
\[
   B^\nu=B[t]=R[x,t]/(\pi-xt),
\]
with $y=t^2$: the standard regular nodal model $\mathcal U$ of
\Cref{ex:nodal-model} in coordinates $(x,t)$.  The normalization map
\[
   \nu:\operatorname{Spec}R[x,t]/(\pi-xt)
   \longrightarrow\mathcal Z,
   \qquad (x,t)\longmapsto(x,\,t^2),
\]
is finite, is an isomorphism over the generic fibre (where $t=\pi/x$ is
already a regular function), and untwists the special fibre: the reduced
SNC fibre $\{xt=0\}$ maps onto $\{x^2y=0\}$, with the component
$\{x=0\}=\operatorname{Spec}\widetilde K[t]$ mapping to the doubled line by
the degree-two map $y=t^2$ on coordinates.  Thus normalization simultaneously
repairs the total space and replaces the nonreduced special fibre by a
semistable one.  Model functions computed on $\mathcal Z$ through
$\nu$-pullback agree with those computed on the normalization; this
invariance is \Cref{lem:model-function-well-defined}.
\end{example}

\begin{corollary}
\label{cor:common-model-comparison}
Quantities attached to a model that are invariant under pullback to a
dominating model may be compared on a common model.
\end{corollary}
\begin{proof}
Apply \Cref{prop:common-domination}. This elementary observation is
the reason that model functions, intersection measures, and nefness tests are
organized as direct or inverse limits over models.
\end{proof}

\section{Flattening and model extension}

Two operations recur later. First, a coherent object defined on the generic
fibre can often be extended after passing to a model. Second, a proper
morphism can often be made flat after a blowup of the base.

\begin{proposition}\label{prop:line-bundle-extension}
Let $K$ be discretely valued, let $X$ be projective over $K$, and let $L$ be a
line bundle on $X$. There exist a projective model $\mathcal X$ of $X$ and
a line bundle $\mathcal L$ on $\mathcal X$ with
$\mathcal L_K\simeq L$.
\end{proposition}
\begin{proof}
Choose an ample line bundle $A$ on $X$.  For sufficiently large $q$, both
$L+qA$ and $qA$ are very ample; fix such $q$ and closed
embeddings
\[
      \iota_1:X\hookrightarrow\mathbf P^{N_1}_K,
      \qquad
      \iota_2:X\hookrightarrow\mathbf P^{N_2}_K
\]
with $\iota_1^*\mathcal O(1)\simeq L+qA$ and
$\iota_2^*\mathcal O(1)\simeq qA$.  Let
$\mathcal X_i\subset\mathbf P^{N_i}_{K^\circ}$ be the flat closure of
$\iota_i(X)$ in the projective space over $K^\circ$
(\Cref{lem:flat-closure-generic-fibre}); each $\mathcal X_i$ is a
projective model of $X$, of finite type because $K^\circ$ is noetherian
here.  By \Cref{prop:common-domination}, applied inside
$\mathcal X_1\times_{K^\circ}\mathcal X_2$, there is a projective model
$\mathcal X$ dominating both, say by
$\rho_i:\mathcal X\to\mathcal X_i$.  Set
\[
   \mathcal L
   =\rho_1^*\mathcal O_{\mathbf P^{N_1}}(1)\big|_{\mathcal X}
   \otimes
   \Bigl(\rho_2^*\mathcal O_{\mathbf P^{N_2}}(1)\big|_{\mathcal X}\Bigr)^{-1}.
\]
On the generic fibre the two pullbacks restrict, through the marked
identifications, to $L+qA$ and $qA$, so
$\mathcal L_K\simeq(L+qA)-(qA)=L$, as required.
\end{proof}

\begin{remark}\label{rem:line-bundle-extension-twist}
The auxiliary twist by $qA$ is a device in the proof: $L$ itself need not
be very ample or effective, but the difference of the two extended very
ample bundles extends $L$ itself.  The $\mathbf Q$-models used for model
metrics serve a different purpose: they permit rational model weights and
rational powers of metrics; they are not needed to overcome an extension
obstruction for $L$.
\end{remark}

The second recurring operation is flattening.  The statement below is the
noetherian model form used in this book; the underlying tag is valid over a
quasi-compact quasi-separated base and has already been exploited in that
generality in \Cref{prop:flat-closure-finite-presentation}.

\begin{theorem}\label{thm:flattening-models}
Let $f:Y\to S$ be a morphism of finite presentation with $S$ noetherian, and
let $\mathcal F$ be a coherent sheaf that is flat over a schematically dense
open set $U\subset S$. There is a blowup $S'\to S$, with center disjoint from
$U$, such that the strict transform of $\mathcal F$ on the strict transform
of $Y$ is flat over $S'$.
\end{theorem}
\begin{proof}
This is precisely flattening by blowup in
\cite[Tag 0815]{stacks-project}.  In the notation of that result, the
schematically dense open is $U$, the center furnished by the theorem is
contained in $S\setminus U$, and its proper transform is the strict
transform used here.  The finite-presentation and noetherian hypotheses in
the present statement are exactly the finiteness hypotheses under which the
result produces one coherent blowup center and a flat strict transform.
Thus the cited theorem gives $S'$ and the asserted flatness.  The different
domination statement \cite[Tag 081T]{stacks-project} is not used.
\end{proof}

\section{SNC, dlt, and altered models}

\begin{definition}\label{def:snc-model}
Assume $K$ is discretely valued. A regular model $\mathcal X$ is an
\emph{SNC model} if its special fibre has simple normal crossings: every
irreducible component $E_i$ of $\mathcal X_s$ is regular, and for every
point of $\mathcal X_s$ the local equations $z_0,\ldots,z_r$ of the
\emph{distinct} components $E_0,\ldots,E_r$ through that point extend to
regular parameters $z_0,\ldots,z_r,w_1,\ldots,w_q$, so that, for a unit
$u$,
\[
                         \pi=u z_0^{b_0}\cdots z_r^{b_r}.
\]
The special fibre is \emph{strictly semistable} if all $b_i=1$.  The
regularity of the components, and the requirement that the $z_i$ cut out
distinct components, are not automatic consequences of the displayed
equation: an $\mathrm I_1$ degeneration of an elliptic curve, whose special
fibre is a single nodal rational curve with local equation $xy=\pi$ at the
node, satisfies the equation at every point with all $b_i=1$, yet its
special fibre is not a simple normal crossings divisor and its dual complex
is a loop rather than a simplicial complex.  This is exactly the case
excluded in \Cref{ex:snc-chain-models} for $n=1$.
\end{definition}

The displayed local equation explains why skeleta are simplicial. A valuation
centered in this chart and monomial in the $z_i$ is determined by nonnegative
numbers $\alpha_i=v(z_i)$ satisfying
\[
                         \sum_{i=0}^r b_i\alpha_i=1.
\]
Thus the local skeleton is a simplex. This construction is carried out
intrinsically in \cref{chap:skeleta}.

\begin{example}[The SNC models of the running nodal degeneration]
\label{ex:snc-chain-models}
\index{semistable model!chain}\index{SNC model!example}\index{degeneration!nodal}
Fix $n\geq1$ and consider the local models
\[
    \mathcal U_n=\operatorname{Spec}A_n,
    \qquad A_n=R[x,y]/(xy-\pi^n),
\]
over a discrete valuation ring $R=K^\circ$, together with their projective
companions $\mathcal C^{(n)}=\{XY=\pi^nZ^2\}\subset\mathbf P^2_R$; for
$n=1$ these are the running nodal model \eqref{eq:running-node} and the
conic \eqref{eq:running-conic}.  We construct all their standard SNC
models.

\emph{The singularity.}  For $n=1$ the total space is regular
(\Cref{ex:nodal-model}).  For $n\geq2$ it is normal---a hypersurface, hence
Serre's condition $S_2$, with singular locus the single point
$x=y=\pi=0$, of codimension two---but not regular: at that point the
relation $xy=\pi^n$ lies in the square of the maximal ideal, so
$\mathfrak m/\mathfrak m^2$ has dimension three while the local ring has
dimension two.  This is the $A_{n-1}$-singularity.  The special fibre is
the reduced union of the two axes for every $n$, so the special fibre alone
does not detect $n$; the total space does.

\emph{One blowup.}  Blow up the closed point $(x,y,\pi)$ of
$\mathcal U_n$, $n\geq2$.  On the $x$-chart, write $y=xy_1$, $\pi=x\pi_1$;
the strict transform of $xy-\pi^n$ is
$y_1=x^{\,n-2}\pi_1^{\,n}$, and eliminating $y_1$ identifies the chart with
\[
     \operatorname{Spec}R[x,\pi_1]/(\pi-x\pi_1),
\]
a regular scheme with a semistable node $x\pi_1=\pi$.  The $y$-chart is
symmetric.  On the $\pi$-chart, write $x=\pi x_1$, $y=\pi y_1$; the strict
transform is
\[
     x_1y_1=\pi^{\,n-2},
\]
that is, a copy of $\mathcal U_{n-2}$.  Thus one point blowup peels off
two regular semistable charts and leaves an $A_{n-3}$-singularity
(a regular point if $n\leq3$).

\emph{The case $n=2$ in full.}  Here the $\pi$-chart is
$x_1y_1=1$, i.e.\ $\operatorname{Spec}R[x_1,x_1^{-1}]$, which is smooth
over $R$ with special fibre $\mathbf G_{m,\widetilde K}$.  The exceptional
curve $E$ is glued from $\{x=0\}$ in the $x$-chart, with coordinate
$\pi_1$, and $\{y=0\}$ in the $y$-chart, with coordinate $x_1=1/\pi_1$; it
is a projective line.  The new special fibre is the chain
$\widetilde D_y\cup E\cup\widetilde D_x$, all multiplicities one, with two
nodes of local equation $uv=\pi$.  Computing divisorial points as in
\Cref{ex:node-divisorial-model-function} places the vertices at
$v(x)=0,1,2$ on the skeleton $[0,2]$: the blowup has subdivided the length-two
edge of $\mathcal U_2$ at its midpoint.  The same model was obtained in
\Cref{ex:vertical-blowup-weil-center} by blowing up the vertical Weil
divisor $D_x$; the two centers $(x,\pi)$ and $(x,y,\pi)$ have the same
normalized blowup here.

\emph{Induction.}  Iterating, after $\lfloor n/2\rfloor$ point blowups all
singularities are resolved and one obtains a regular strictly semistable
model $\widetilde{\mathcal U}_n$ whose special fibre is a chain
\[
     \widetilde D_y=E_0,\ E_1,\ \ldots,\ E_{n-1},\ E_n=\widetilde D_x,
\]
with $E_i$ meeting $E_{i+1}$ in a single node of local equation
$uv=\pi$ and all multiplicities equal to one.  The divisorial point of
$E_i$ is the monomial valuation with $v(x)=i$, $v(y)=n-i$.  The dual
complex is therefore the segment $[0,n]$ subdivided into unit edges: the
model-theoretic incarnation of the rule from
\cref{sec:running-semistable} that a node $xy=a$ contributes an edge of
length $v(a)$.  The intermediate models realize the partial subdivisions
with vertices $\{0,1,\ldots,k\}\cup\{n-k,\ldots,n\}$, and blowing up nodes
of any semistable model further subdivides edges without changing the path
metric; this is the cofinal system of subdivisions taken up in
\cref{sec:skeleton-refinement}.

\emph{Globalization.}  For the projective family $\mathcal C^{(n)}$ the
computation is identical in the affine chart $Z=1$ and trivial elsewhere,
so the special fibre of the resolved model
$\widetilde{\mathcal C}^{(n)}$ is a chain of $n+1$ projective lines: the
strict transforms of the two lines $\{X=0\}$ and $\{Y=0\}$ at the ends and
$n-1$ exceptional curves between them.  For a projective degeneration of
curves with smooth generic fibre and several nodes, each node carries out
this computation independently; in particular, for a nodal degeneration of
an \emph{elliptic} curve whose special fibre is an irreducible rational
curve with one node of local equation $xy=\pi^{\,n}$, with $n\ge2$, the resolved SNC
model has a cycle of $n$ rational curves as special fibre and a circle of
circumference $n$ as dual complex.  This is the source of the nontrivial
essential skeleta discussed in \cref{sec:essential-skeleton}.
For $n=1$ the irreducible nodal component is not an SNC divisor in the
strict sense used here; an additional blowup is needed, and its exceptional
component occurs with multiplicity two.
\end{example}

A \emph{dlt model} is a normal model for which the pair consisting of the
model and the reduced special fibre has divisorial log terminal singularities.
Its strata still yield a dual complex, while dlt modifications are frequently
available when regular SNC models are not. We defer the minimal-model
definitions to the chapter where they are needed.

\paragraph{Model-availability input ledger.}
The following inputs will be used with the indicated scope.
\begin{enumerate}
\item In the equicharacteristic-zero, finite-type situations used by the
  original BFJ theory, resolution and embedded resolution provide regular
  SNC models after suitable modifications.
\item In arbitrary characteristic, alteration theorems often provide, after
  finite extension and modification, a proper generically finite cover with
  regular total space and normal-crossings boundary. The precise conclusion
  depends on the alteration theorem invoked, and the cover is not in general
  birational.
\item In mixed or positive characteristic, the existence of a regular SNC
  model in arbitrary dimension is not assumed.
\end{enumerate}

Part (1) is the resolution and embedded-resolution theorem
\cite{Hironaka1964I,Hironaka1964II} in its model form; part (2) is the
semistable alteration theorem \cite{deJong1996}. These are deep external
inputs. Part (3) is a restriction on our later arguments. In particular, an
alteration cannot simply replace a resolution in an intersection computation:
degrees, descent, and possible inseparability must be tracked.

\section{Intersection theory on a model}

In this section $K^\circ$ is a discrete valuation ring, $\mathcal X$ is a
proper flat model of relative dimension $n$, and all schemes are noetherian.
If $D=\sum_Ea_EE$ is a vertical Cartier divisor and
$\mathcal L_1,\ldots,\mathcal L_n$ are line bundles on $\mathcal X$, define
\[
 (D\cdot\mathcal L_1\cdots\mathcal L_n)_{\mathcal X}
   :=\sum_E a_E
      \deg_E\bigl(c_1(\mathcal L_1|_E)\cdots
                  c_1(\mathcal L_n|_E)\cap[E]\bigr).
\]
Residue-field extension degrees are included in $\deg_E$. Notice that the
vertical divisor, not merely its associated line bundle, is part of this
relative intersection datum.

\begin{proposition}\label{prop:model-intersection-basic}
The vertical intersection number is linear in $D$, symmetric and multilinear
in $\mathcal L_1,\ldots,\mathcal L_n$, and depends only on their restrictions
to the support of $D$.
\end{proposition}
\begin{proof}
Linearity in $D$ is linearity of the associated cycle $[D]$.  Symmetry and
multilinearity in the line bundles are the corresponding properties of first
Chern class operators on the Chow groups of the proper
$\widetilde K$-schemes $E$.  Since every summand is computed on a component
of $|D|$, replacing a line bundle by another with the same restriction to
$|D|$ leaves every summand unchanged.
\end{proof}

Before proving the projection formula we record the intersection-theoretic
framework precisely.  Cycles, proper pushforward, flat pullback, and
intersection with Cartier divisors and Chern classes of line bundles are
available, with all the identities of \cite[Chapters~1--2]{Ful}, for
schemes of finite type over the regular one-dimensional base
$\operatorname{Spec}K^\circ$; this extension is
\cite[\S20.1]{Ful}.  For a zero-cycle $\alpha=\sum_i n_i[p_i]$ on a scheme
proper over $\widetilde K$ we write
$\deg\alpha=\sum_in_i[\kappa(p_i):\widetilde K]$; the degree is computed by
pushing forward to $\operatorname{Spec}\widetilde K$ and is therefore
unchanged under proper pushforward between proper
$\widetilde K$-schemes \cite[\S1.4]{Ful}.  All classes below are supported
on special fibres, which are proper over $\widetilde K$ because the models
are proper over $K^\circ$, so degrees are always defined.

\begin{lemma}
\label{lem:vertical-intersection-degree}
Let $D=\sum_Ea_EE$ be a vertical Cartier divisor on $\mathcal X$ and let
$[D]=\sum_Ea_E[E]\in Z_n(\mathcal X_s)$ be its Weil cycle.  Then
\[
 (D\cdot\mathcal L_1\cdots\mathcal L_n)_{\mathcal X}
 =\deg\bigl(c_1(\mathcal L_1)\cdots c_1(\mathcal L_n)\cap[D]\bigr).
\]
Moreover, if $\mathcal X$ is integral, then $[D]$ represents the
Cartier-divisor intersection with support
\[
                      D\mathbin{\cdot}[\mathcal X]\in A_n(|D|).
\]
For the closed immersion $i:|D|\hookrightarrow\mathcal X$, its pushforward
is
\[
 i_*\bigl(D\mathbin{\cdot}[\mathcal X]\bigr)
       =c_1(\cO(D))\cap[\mathcal X]\qquad\text{in }A_n(\mathcal X),
\]
but the supported class and its pushforward must not be identified.
\end{lemma}
\begin{proof}
Both sides of the displayed identity are additive in the cycle $[D]$, so it
suffices to treat a single prime component $E$, where it is the definition
of $\deg_E$.  The supported Cartier-divisor intersection
$D\mathbin{\cdot}[\mathcal X]$ is represented by the associated Weil cycle
$[D]$ \cite[\S2.3]{Ful}.  Its pushforward to $\mathcal X$ is the ordinary
first-Chern-class action by the same construction, giving the last displayed
identity.
\end{proof}

\begin{theorem}\label{thm:model-projection-formula}
Let $\mathcal X$ and $\mathcal Y$ be integral proper flat models of the same
relative dimension, and let
$\rho:\mathcal Y\to\mathcal X$ be a proper dominant morphism whose generic
fibre is generically finite of degree $d$. For every vertical Cartier divisor
$D$ on $\mathcal X$,
\[
 (\rho^*D\cdot\rho^*\mathcal L_1\cdots\rho^*\mathcal L_n)_{\mathcal Y}
       =d\,(D\cdot\mathcal L_1\cdots\mathcal L_n)_{\mathcal X}.
\]
In particular the number is unchanged under a dominating birational model.
\end{theorem}
\begin{proof}
\emph{Step 1: $\rho^*D$ is a vertical Cartier divisor and
$\rho_*[\mathcal Y]=d[\mathcal X]$.}
Pullback of Cartier divisors along the dominant morphism $\rho$ is defined
by pulling back local equations, and
$\operatorname{Supp}\rho^*D\subset\rho^{-1}(\operatorname{Supp}D)
\subset\rho^{-1}(\mathcal X_s)=\mathcal Y_s$,
since $\rho$ is a $K^\circ$-morphism.  For the second assertion, proper
pushforward of the fundamental cycle is governed by the function-field
degree $[R(\mathcal Y):R(\mathcal X)]$ \cite[\S1.4]{Ful}.  By
\Cref{prop:vertical-horizontal}\,(3) the generic fibres are dense in both
models, so $R(\mathcal Y)=K(\mathcal Y_K)$ and
$R(\mathcal X)=K(X)$, and the extension degree is the degree $d$ of the
generically finite generic fibre of $\rho$.  Hence
$\rho_*[\mathcal Y]=d[\mathcal X]$.

\emph{Step 2: projection for the supported divisor class.}  Write
$\rho_D:|\rho^*D|\to|D|$ for the proper morphism induced by $\rho$.
The projection formula for intersection with a Cartier divisor, formulated
with support, gives
\[
 \rho_{D*}\bigl(\rho^*D\mathbin{\cdot}[\mathcal Y]\bigr)
   =D\mathbin{\cdot}\rho_*[\mathcal Y]
   =d\bigl(D\mathbin{\cdot}[\mathcal X]\bigr)
       \qquad\text{in }A_n(|D|);
\]
see \cite[Proposition~2.3(c)]{Ful}.  This is the essential supported
identity.  Replacing it by an identity involving only
$c_1(\cO(D))\cap[\mathcal X]$ would lose information whenever that class
vanishes after pushforward to the total space.

\emph{Step 3: apply the remaining Chern classes.}  The projection formula
for first Chern class operators \cite[Proposition~2.5(c)]{Ful}, applied on
the supports once for each $\mathcal L_i$, now yields
\[
 \begin{aligned}
 &\rho_{D*}\!\left(
 c_1(\rho^*\mathcal L_1)\cdots c_1(\rho^*\mathcal L_n)
 \cap\bigl(\rho^*D\mathbin{\cdot}[\mathcal Y]\bigr)\right)\\
 &\qquad=d\,c_1(\mathcal L_1)\cdots c_1(\mathcal L_n)
 \cap\bigl(D\mathbin{\cdot}[\mathcal X]\bigr).
 \end{aligned}
\]
Both sides are zero-cycle classes supported on the special fibres, and
$\rho_D$ is a proper morphism between closed subschemes of proper
$\widetilde K$-schemes.

\emph{Step 4: take degrees.}  Degrees of zero-cycles supported on the
special fibres are preserved by $\rho_{D*}$, since pushing forward to
$\operatorname{Spec}\widetilde K$ factors through $\rho_D$.  Taking
$\deg$ of the identity in Step 3 and using
\Cref{lem:vertical-intersection-degree} on both models yields the theorem.
The final sentence is the case $d=1$.
\end{proof}

\begin{proposition}\label{prop:special-fibre-degree-conservation}
Assume in addition that $\mathcal X$ is normal.
Writing $\mathcal X_s=\sum_Eb_EE$, one has
\[
 \sum_Eb_E(\mathcal L_1|_E\cdots\mathcal L_n|_E)_E
   = (L_1\cdots L_n)_X,
\]
where $L_i=\mathcal L_i|_X$.
\end{proposition}
\begin{proof}
We spell out the Euler-characteristic argument, which is conservation of
degree in a proper flat family.  Two external inputs are used, each cited
exactly.

\emph{Step 1: multiplicities as lengths.}  The coefficient $b_E$ was
defined in the normal case as $\operatorname{ord}_E(\pi)$.  We first check
that it equals the length of the scheme-theoretic special fibre at the
generic point $\eta_E$ of $E$:
\[
 b_E=\operatorname{length}_{\mathcal O_{\mathcal X,\eta_E}}
     \bigl(\mathcal O_{\mathcal X,\eta_E}/(\pi)\bigr).
\]
Indeed, on the normal model the local ring $\mathcal O_{\mathcal
X,\eta_E}$ is a discrete valuation ring with valuation
$\operatorname{ord}_E$, and for an element of order $b$ the quotient by it
has length $b$.  Consequently the fundamental cycle of the scheme
$\mathcal X_s$ is $\sum_Eb_E[E]$ plus nothing in dimension $n$: any other
component of $\mathcal X_s$ has dimension $<n$.

\emph{Step 2: invariance of Euler characteristics.}  For integers
$m_1,\ldots,m_n\geq0$ put
$\mathcal M(\mathbf m)=\mathcal L_1^{\otimes m_1}\otimes\cdots\otimes
\mathcal L_n^{\otimes m_n}$.  Since $\mathcal X$ is proper and flat over
$K^\circ$ and $\mathcal M(\mathbf m)$ is a line bundle, hence flat, the
Euler characteristic of its fibres is constant in the family
\cite[7.9.4]{EGAIII2}:
\[
   \chi\bigl(X,\mathcal M(\mathbf m)|_X\bigr)
   =\chi\bigl(\mathcal X_s,\mathcal M(\mathbf m)|_{\mathcal X_s}\bigr)
   \qquad\text{for all }\mathbf m .
\]

\emph{Step 3: Snapper polynomials on both fibres.}  For a proper scheme
$Y$ over a field $k$, a coherent sheaf $F$, and line bundles
$M_1,\ldots,M_n$, the function
$\mathbf m\mapsto\chi(Y,F\otimes M_1^{m_1}\otimes\cdots\otimes M_n^{m_n})$
is a polynomial of total degree at most $\dim\operatorname{Supp}F$, and
when $\dim\operatorname{Supp}F=n$ the coefficient of $m_1\cdots m_n$ in it
is the mixed intersection number
$\deg(c_1(M_1)\cdots c_1(M_n)\cap[F]_n)$ against the $n$-dimensional part
of the cycle of $F$; this is Snapper's theorem, in the form of asymptotic
Riemann--Roch \cite[Example~18.3.6]{Ful}, applied to
$M_1^{m_1}\otimes\cdots\otimes M_n^{m_n}$ and polarized.

On the generic fibre, with $F=\mathcal O_X$ and $[X]$ the fundamental
cycle, the coefficient of $m_1\cdots m_n$ is $(L_1\cdots L_n)_X$.  On the
special fibre, with $F=\mathcal O_{\mathcal X_s}$, Step 1 gives
$[\mathcal X_s]_n=\sum_Eb_E[E]$, so the same coefficient is
$\sum_Eb_E(\mathcal L_1|_E\cdots\mathcal L_n|_E)_E$; components of
$\mathcal X_s$ of dimension $<n$ contribute polynomials of total degree
$<n$, which have no $m_1\cdots m_n$ term.  By Step 2 the two polynomials
are equal, hence so are the two coefficients.
\end{proof}

\begin{example}[Degree conservation on the running conic]
\label{ex:degree-conservation-conic}
For the conic model $\mathcal C=\{XY=\pi Z^2\}$ of
\eqref{eq:running-conic}, with
$\mathcal L=\mathcal O_{\mathbf P^2_R}(1)|_{\mathcal C}$, the special fibre
is $C_0+C_1$, two lines with $b_{C_i}=1$ and
$\deg\mathcal L|_{C_i}=1$, while $\deg_CL=2$; the identity reads
$1+1=2$ and is the degree computation behind the Chambert--Loir measure of
\Cref{prop:running-conic-measure}.  For the semistable chain model
$\widetilde{\mathcal C}^{(2)}$ of \Cref{ex:snc-chain-models}, the special
fibre is a chain of three lines of degrees $1,0,1$ under the pullback of
$\mathcal L$, and the identity reads $1+0+1=2$: domination redistributes
degree along the chain but conserves the total.
\end{example}

\begin{remark}\label{rem:principal-special-fibre-relative-degree}
Let $i:\mathcal X_s\hookrightarrow\mathcal X$.  The supported Cartier-divisor
class
\[
              \mathcal X_s\mathbin{\cdot}[\mathcal X]
                    =[\mathcal X_s]\in A_n(\mathcal X_s)
\]
can be nonzero even though its pushforward is zero: since
$\mathcal X_s=\operatorname{div}(\pi)$ is principal,
$\cO(\mathcal X_s)\simeq\cO_{\mathcal X}$ and hence
\[
 i_*[\mathcal X_s]
   =c_1(\cO(\mathcal X_s))\cap[\mathcal X]=0
       \qquad\text{in }A_n(\mathcal X).
\]
The relative vertical intersection takes the degree only after the other
Chern class operators have acted on the supported class in the proper special
fibre.  It is therefore generally nonzero and, by
\Cref{prop:special-fibre-degree-conservation}, equals the generic degree.
Confusing the supported class with its zero pushforward would incorrectly
force every Chambert--Loir mass to vanish.
\end{remark}

\section{Relative nefness and model semipositivity}

We retain the discretely valued hypothesis of the preceding section for
these intersection-theoretic nefness tests.

\begin{definition}\label{def:relative-nefness-model-semipositivity}
A line bundle $\mathcal L$ on a proper model $\mathcal X$ is
\emph{relatively nef} if
\[
                            (\mathcal L\cdot C)\geq0
\]
for every integral proper curve $C$ contained in a fibre of
$\mathcal X\to\operatorname{Spec}K^\circ$. A $\mathbf Q$-model of $L$ is
\emph{semipositive} if its model line bundle is relatively nef.
\end{definition}

Relative nefness can be checked on the special fibre. Indeed, the flat
closure of an integral proper curve in $X$ is a proper flat family over
$K^\circ$. The Euler-characteristic argument used in
\cref{prop:special-fibre-degree-conservation} expresses its generic degree
as the sum of the degrees on the special-fibre components, with their
positive length multiplicities \cite[7.9.4]{EGAIII2}. Nonnegative degrees
on all special-fibre curves therefore imply nonnegative degree on the
generic curve as well. Relative nefness is preserved by pullback.

\begin{proposition}
\label{prop:nef-pullback}
Let $\rho:\mathcal Y\to\mathcal X$ be a vertical modification. If
$\mathcal L$ is relatively nef on $\mathcal X$, then $\rho^*\mathcal L$ is
relatively nef on $\mathcal Y$.
\end{proposition}
\begin{proof}
Let $C\subset\mathcal Y_s$ be an integral curve. If $\rho(C)$ is a point, the
restriction of $\rho^*\mathcal L$ to $C$ is pulled back from a one-dimensional
vector space, hence has degree zero. If $\rho(C)$ is a curve, the projection
formula gives
\[
                 (\rho^*\mathcal L\cdot C)
       =\deg(C/\rho(C))(\mathcal L\cdot\rho(C))\geq0.
\]
\end{proof}

A model metric is semipositive when it admits a relatively nef model
representative. The proposition shows that every domination of this
representative remains nef. The comparison among algebraic, formal, and piecewise linear semipositive metrics
is proved in \cite{GM19} and developed in \cref{chap:metrics}.

\begin{example}[Relative nefness on the chain model]
\label{ex:chain-nef-functions}
\index{nef!relatively nef example}\index{model function!nefness}
We compute the nef region among vertical twists on the semistable chain
model $\widetilde{\mathcal C}=\widetilde{\mathcal C}^{(2)}$ of
\Cref{ex:snc-chain-models}, obtained from
$\mathcal C^{(2)}=\{XY=\pi^2Z^2\}$ by one blowup.  Its special fibre is the
chain $C_0,E,C_2$ of projective lines, all multiplicities one, with
divisorial points at $v(x)=0,1,2$ on the skeleton $[0,2]$.  Write
$\rho:\widetilde{\mathcal C}\to\mathcal C^{(2)}$ and
$\mathcal L_0=\rho^*\bigl(\mathcal O_{\mathbf P^2_R}(1)|_{\mathcal
C^{(2)}}\bigr)$.

First the intersection matrix.  The fibre degrees of $\mathcal L_0$ are
$(\mathcal L_0\cdot C_0)=1$, $(\mathcal L_0\cdot E)=0$
(the curve $E$ is contracted by $\rho$, so $\mathcal L_0|_E$ is trivial),
and $(\mathcal L_0\cdot C_2)=1$.  The chain incidences are
$(C_0\cdot E)=(E\cdot C_2)=1$ and $(C_0\cdot C_2)=0$.  Since the full
fibre $C_0+E+C_2=\operatorname{div}(\pi)$ has intersection zero with every
fibre curve, the self-intersections follow:
\[
   (C_0^2)=(C_2^2)=-1,\qquad (E^2)=-2 .
\]

Now let $D=a_0C_0+a_1E+a_2C_2$ be a vertical divisor with $a_i\in\mathbf Z$,
or with $a_i\in\mathbf Q$ after passing to $\mathbf Q$-models. The same
inequalities describe the corresponding region in the space of real
vertical classes. The model function takes the values $a_0,a_1,a_2$ at the
three divisorial points. Set
$\mathcal L=\mathcal L_0\otimes\cO(D)$, a model of $\mathcal O(1)|_C$ on
which the metric has been perturbed by $\varphi_D$.  Testing against the
three fibre curves:
\[
 \begin{aligned}
 (\mathcal L\cdot C_0)&=1-a_0+a_1,\\
 (\mathcal L\cdot E)&=a_0-2a_1+a_2,\\
 (\mathcal L\cdot C_2)&=1+a_1-a_2 .
 \end{aligned}
\]
Thus $\mathcal L$ is relatively nef if and only if the function
$i\mapsto a_i$ on the subdivided skeleton is convex at the interior vertex
and its outgoing slopes at the two endpoints are bounded below by $-1$.
These are exactly the conditions under which the piecewise affine metric
weight on the segment is semipositive for the degree-two polarization, in
accordance with the curvature computation
\eqref{eq:running-perturbed-measure} and the nef interval
\eqref{eq:running-nef-interval} of the length-one model. For a consistency
check, the sum of the three numbers is
$2$ for every $(a_0,a_1,a_2)$, as
\Cref{prop:special-fibre-degree-conservation}
requires. Adding a multiple of the full fibre changes the metric by a
constant and leaves each number unchanged. Taking $D=E$ alone gives
$(\mathcal L\cdot E)=-2<0$, so the exceptional curve itself is never a nef
direction.  The three intersection numbers will be reinterpreted in
\cref{chap:metrics} as the masses of the Chambert--Loir measure at the
three divisorial points.
\end{example}

\section{The inverse system of models}

Let $\mathfrak M_X$ be the category of proper models of $X$, with morphisms
given by domination. \Cref{prop:common-domination} says that it is
cofiltered after the appropriate finiteness restrictions. Two constructions
will repeatedly occur:
\[
       \varinjlim_{\mathcal X\in\mathfrak M_X}
       \operatorname{Div}_0(\mathcal X)_{\mathbf Q},
       \qquad
       \varprojlim_{\mathcal X\in\mathfrak M_X^{\mathrm{sk}}}
       |\Delta(\mathcal X)|.
\]
The first uses pullback of vertical Cartier divisors and produces Cartier
b-divisors, or equivalently rational model functions.  In the second,
$\mathfrak M_X^{\mathrm{sk}}$ denotes a specified cofiltered system of models
that actually carry skeleta---for example SNC models in the
equicharacteristic-zero regime.  It uses the retractions between their dual
complexes and helps recover the normalized valuative space.  We do not write
$\Delta(\mathcal X)$ for an arbitrary singular model.  The transition maps
and the class of allowed models are part of the definitions; neither limit
is a merely set-theoretic union.

Both limits are only as good as the index category, and this chapter has
now supplied everything the index category needs: cofilteredness is
\Cref{prop:common-domination}, sharpened by
\Cref{prop:flat-closure-finite-presentation} over nonnoetherian valuation
rings; the transition maps on divisors are pullbacks, compatible by
\Cref{lem:model-function-well-defined}; and the transition maps on dual
complexes come from the subdivision behaviour seen concretely in
\Cref{ex:snc-chain-models}.  The formalism of Cartier and Weil b-divisors
built on the first limit is developed in \cref{app:intersection-bdivisors},
beginning with \Cref{def:appF-cartier-bdivisor}; the second limit is the
subject of \cref{chap:skeleta}.

\begin{proposition}
\label{prop:divisor-model-function-injective-mod-fibre}
Assume $X$ is proper and geometrically connected and $K$ is discretely
valued. On a fixed normal model $\mathcal X$, the map
\[
 D\longmapsto\varphi_D
\]
from vertical $\mathbf Q$-Cartier divisors to model functions is injective.
Moreover $\varphi_{\mathcal X_s}=1$, so adding
$c\mathcal X_s$ adds the constant $c$.
\end{proposition}
\begin{proof}
At the divisorial point $x_E$,
$\varphi_D(x_E)=a_E/b_E$. If $\varphi_D=0$, all coefficients $a_E$ vanish,
so $D=0$. The formula for $\mathcal X_s$ follows because its coefficient
along $E$ is $b_E$. Linearity gives the final assertion.
\end{proof}

\section{A dictionary of the model language}
\label{sec:models-dictionary}
\index{dictionary!models}

The chapter has built one half of several correspondences whose other
halves appear in \cref{chap:berkovich,chap:skeleta,chap:metrics}.  We
record the dictionary in the state it has reached, with pointers to where
each entry is established or completed. Entries involving $b_E$ or $\pi$
use the discrete normalization fixed above; the component-valuation formula
also assumes a normal model.

\medskip
\noindent
\begin{tabularx}{\textwidth}{@{}>{\raggedright\arraybackslash}p{.25\textwidth}%
  >{\raggedright\arraybackslash}X%
  >{\raggedright\arraybackslash}p{.25\textwidth}@{}}
\toprule
model datum & valuative or analytic datum & where\\
\midrule
component $E\subset\mathcal X_s$, multiplicity $b_E$
 & divisorial point $x_E$, $v_E=\tfrac1{b_E}\operatorname{ord}_E$
 & \Cref{prop:component-valuation}\\
vertical Cartier divisor $D$
 & model function $\varphi_D=-\log|f_D|$
 & \Cref{lem:model-function-well-defined}\\
effective vertical divisor
 & nonnegative model function
 & sign convention, \Cref{ex:node-divisorial-model-function}\\
special fibre $\mathcal X_s=\operatorname{div}(\pi)$
 & constant function $1$
 & \Cref{prop:divisor-model-function-injective-mod-fibre}\\
domination $\mathcal Y\to\mathcal X$
 & pullback compatibility; new divisorial points
 & \Cref{prop:common-domination}\\
blowup of a semistable node
 & subdivision of a skeleton edge
 & \Cref{ex:snc-chain-models}, \cref{sec:skeleton-refinement}\\
SNC model
 & simplicial skeleton, PA functions
 & \cref{sec:skeleton-snc,sec:skeleton-model-functions}\\
vertical twist on the running chain model
 & relative nefness as convexity and slope inequalities
 & \Cref{ex:chain-nef-functions}\\
relatively nef model bundle
 & semipositive model metric
 & \cref{chap:metrics}, \cite{GM19}\\
vertical intersection numbers
 & Chambert--Loir masses
 & \cref{chap:metrics}, \cref{chap:model-calculations}\\
inverse system of models
 & Berkovich space; Cartier b-divisors
 & \cref{chap:berkovich}, \cref{app:intersection-bdivisors}\\
\bottomrule
\end{tabularx}
\medskip

Each right-hand entry that lives on the analytic side is, strictly
speaking, defined only in the later chapter cited; the table asserts that
the model-side constructions of the present chapter are exactly the
combinatorial shadows that those chapters will attach to them, and no
entry is used before its defining chapter.

\section{What is cited and what is proved}

The following ledger will govern later chapters.
\begin{enumerate}
\item Elementary valuation-ring algebra, component valuations, common graph
  closures, vertical blowups, and the intersection identities above may be
  used directly by chapter and proposition number.
\item Raynaud generic fibres and invariance under admissible formal blowup are
  cited to formal analytic geometry over nontrivially valued fields.
\item Flattening is cited to the precise Stacks Project tags. The
  noetherian model version is distinguished from the general
  finite-presentation argument over a valuation ring.
\item Resolution, semistable reduction, and alterations are named external
  inputs; their characteristic and degree hypotheses remain visible.
\item Chow intersection theory is used only for noetherian proper models over
  a discrete valuation ring unless a later chapter explicitly supplies a
  more general theory.
\item Over a valuation ring with dense value group, the
  torsion-freeness, flat-closure, and finite-presentation statements
  (\Cref{lem:torsion-free-flat-valuation},
  \Cref{lem:flat-closure-generic-fibre},
  \Cref{prop:flat-closure-finite-presentation}) may be invoked; every
  statement mentioning a uniformizer, multiplicities, or intersection
  numbers is restricted to the discretely valued case, as itemized in
  \Cref{ex:model-dense-value-group}.
\end{enumerate}

\section{Rank-one topology and higher-rank integral data}
\label{sec:models-higher-rank-integral-data}
\index{valuation ring!higher-rank integral data}

The higher-rank warning in
\Cref{rem:vertical-horizontal-higher-rank} has an analytic counterpart.
Here assume the absolute value is nontrivial. One can keep the complete
real-valued field $K$ and its norm topology, while choosing an additional
valuation ring
\[
 K^{\circ\circ}\subset K^+\subset K^\circ,\qquad
 k=\widetilde K,\qquad V=K^+/K^{\circ\circ}\subset k.
\]
Here $K^+$ may have higher rank. The adic base
$\operatorname{Spa}(K,K^+)$ records that integral data, whereas
$\operatorname{Spa}(K,K^\circ)$ is its rank-one locus.
The recent account \cite{KaibWerner26HigherRankV1} studies this situation.
The real absolute value remains the one fixed at the start of the chapter;
the extra integral directions are not new real-valued seminorm coordinates.

For a concrete complete field, take $K=\mathbf Q_p((t))$ with
$|t|=\rme^{-1}$ and every nonzero constant of absolute value one.
For $f\ne0$, write $f=t^n(c_n+c_{n+1}t+\cdots)$ with
$c_n\ne0$. Composing the $t$-order with the normalized $p$-order on
the residue field gives
\begin{equation}
 \begin{aligned}
 w(f)&=(n,\operatorname{ord}_p(c_n))\in\mathbf Z^2_{\rm lex},\\
 K^+&=\mathbf Z_p+t\mathbf Q_p[[t]],\\
 K^\circ&=\mathbf Q_p[[t]],\qquad V=\mathbf Z_p.
 \end{aligned}
 \tag{3.1}\label{eq:models-composed-valuation}
\end{equation}
Set $w(0)=+\infty$. Leading coefficients multiply, so $w(fg)=w(f)+w(g)$.
For a sum, a lower $t$-order dominates; at equal $t$-order the
$p$-adic inequality controls the leading coefficient, and its cancellation
raises the first coordinate. Thus this is an ordered-group valuation.
The integral-membership rule is $n>0$, or $n=0$ with
$\operatorname{ord}_p(c_n)\geq0$, giving the displayed ring.
In particular $w(p)=(0,1)$ and $w(t/p^N)=(1,-N)$.

The prime chain is
\[
 (0)\subsetneq\mathfrak p:=t\mathbf Q_p[[t]]
 \subsetneq\mathfrak m:=p\mathbf Z_p+t\mathbf Q_p[[t]].
\]
The adic base has the two nongeneric points, and its rank-one point
corresponds to $\mathfrak p$. The ideal $\mathfrak p$ is not finitely
generated: any finite nonzero set of generators has a least
lexicographic value, and some $t/p^N$ has smaller positive lexicographic
value. Divisibility as in \Cref{lem:valuation-ideals-principal}
then excludes generation. This does not contradict the discreteness of
the \emph{topological} $t$-valuation. Nor can its full lexicographic
refinement embed as an ordered subgroup of $\mathbf R$: a positive
image of $(1,0)$ would exceed every integer multiple of the positive
image of $(0,1)$. Taking the first coordinate is a coarsening that
forgets the residue-field direction.

The integral choice matters even when proper analytic objects can be
recovered. Proposition~3.12 of \cite{KaibWerner26HigherRankV1} gives an
equivalence for proper finite-type adic spaces under passage to the
rank-one base, with the theorem's relative base assumptions. Its
partially proper full-faithfulness input is Proposition~3.3.
Admissible formal models have an additional gluing datum: Proposition~3.10
combines a rank-one formal model with a compatible flat, locally finitely
presented scheme over $V$. Smoothness and properness of the reduction
are tested over $\operatorname{Spec}V$, not just over its generic field $k$.
Consequently proper analytic recovery does not assert that every choice
of smooth integral reduction is recovered uniquely.

An elliptic example makes the last distinction visible without a new
uniformization theorem. For a prime $p\geq5$, put
$a=(1728p-1)^{-1}$ and consider
\[
 E:\ y^2=x^3+3ax+2a,
 \qquad
 \Delta=-\frac{1728^2p}{(1728p-1)^3},\qquad j(E)=\frac1p.
\]
These invariants follow from $\Delta=-16(4A^3+27B^2)$ and
$c_4=-48A$ for $A=3a$, $B=2a$, with $j=c_4^3/\Delta$.
The nonzero constant $\Delta$ is a unit of $K^\circ$, so the equation
gives good $t$-adic reduction. But $w(j)=(0,-1)$, and $j\notin K^+$.
The standard necessary integrality of $j$ for good elliptic reduction
excludes an elliptic scheme over $K^+$ with this generic fiber.
Equivalently, such a scheme would reduce modulo $\mathfrak p$ to an
elliptic scheme over $\mathbf Z_p$ whose generic fiber has nonintegral
$j$. This uses the necessary good-reduction criterion; it does not
prove a general converse. The algebraically closed classification in
Example~6.14 of the source has stronger hypotheses than this worked field.

For abeloids, the source's Raynaud extension has a torus and a
\emph{generically good} abeloid quotient after a finite separable
affinoid extension. Good reduction over all of $K^+$ is a stronger
condition; the pushout comparison in its Corollary~6.9 distinguishes them.
Full algebraic N\'eron models and the two marked formal variants also
have different mapping domains. The exact interfaces and a component-group
warning are recorded in \cref{sec:appL-higher-rank-abeloid-boundary}.
These attributed statements use external adic and formal-model foundations.
They do not extend the book's real-valued Monge--Amp\`ere or energy
constructions to an arbitrary higher-rank ordered value group.

\section*{Source notes}

The model language follows Berkovich and Raynaud as presented in
\cite{Berk12,BoschLutkebohmert1993}. Model functions, divisorial points, and
their sign convention follow \cite[\S2]{BFJ16}. Semipositive model metrics
and the comparison of formal, algebraic, and piecewise linear terminology are
treated in \cite{GM19}. For flattening and normalization we cite stable tags
of \cite{stacks-project}. Intersection theory is from \cite{Ful};
alterations from \cite{deJong1996}.

The valuation-ring algebra of the first section is classical and goes back
to Krull; systematic references are
\cite[Chapter~VI]{ZariskiSamuelII} and the valuation-ring sections of
\cite{stacks-project}.  Flattening by blowup, used here both in its
noetherian model form (\Cref{thm:flattening-models}) and over an arbitrary
valuation ring (\Cref{prop:flat-closure-finite-presentation}), originates
with Raynaud--Gruson \cite{RaynaudGruson1971}; the fact that flat
finite-type modules over a domain are finitely presented is also theirs.
The Euler-characteristic invariance in flat proper families is
\cite[7.9.4]{EGAIII2}, and the excellence of complete discrete valuation
rings, which drives every finiteness-of-normalization statement in the
chapter, is part of \cite[Scholie~7.8.3]{EGAIV2}.  Intersection theory
over a discrete valuation ring, rather than a field, is justified by
\cite[\S20.1]{Ful}; a full-strength intersection theory on models over
arbitrary rank-one valuation rings, which this chapter deliberately does
not assume, is developed through formal geometry and delta-forms in
\cite{GK17,GK19}, and the semistable chain computations of
\Cref{ex:snc-chain-models} are special cases of the semistable reduction
circle of ideas for curves; in arbitrary dimension and characteristic the
corresponding statement is the alteration theorem of \cite{deJong1996},
whose exact scope is recorded in the model-availability ledger above and
in \cref{sec:appB-alterations}.  The appendix
\cref{app:stacks} lists every Stacks Project tag used here with its
hypotheses and first point of application; the appendix
\cref{app:normalizations} records the sign conventions that
\Cref{ex:node-divisorial-model-function} tests on the running example.
